\documentclass[11pt,a4paper,reqno]{amsart}

\usepackage[utf8]{inputenc}
\usepackage[T1]{fontenc}
\usepackage{lmodern}
\usepackage{amsmath,amssymb,amsfonts,amsthm,mathtools}
\usepackage{mathrsfs}
\usepackage[expansion=false]{microtype}
\usepackage{booktabs}
\usepackage{tabularx}
\usepackage[dvipsnames]{xcolor}
\usepackage{hyperref}
\usepackage{geometry}
\usepackage{enumitem}

\geometry{
  top=3cm,
  bottom=3cm,
  left=2.8cm,
  right=2.8cm,
  headheight=14pt
}

\hypersetup{
  colorlinks=true,
  linkcolor=NavyBlue,
  citecolor=Mahogany,
  urlcolor=TealBlue,
  pdfauthor={Reinaldo M. Silva-Filho},
  pdftitle={Beyond the Spectrum II: Metric Measure Geometry, Persistent Homology, and Non-Commutative Invariants of Functional Tensor Manifolds}
}

% --- Theorem Environments ---
\newtheorem{theorem}{Theorem}[section]
\newtheorem{lemma}[theorem]{Lemma}
\newtheorem{proposition}[theorem]{Proposition}
\newtheorem{corollary}[theorem]{Corollary}

\theoremstyle{definition}
\newtheorem{definition}[theorem]{Definition}
\newtheorem{example}[theorem]{Example}
\newtheorem{remark}[theorem]{Remark}
\newtheorem{construction}[theorem]{Construction}

\numberwithin{equation}{section}
\raggedbottom

% --- Custom Notations ---
\newcommand{\R}{\mathbb{R}}
\newcommand{\C}{\mathbb{C}}
\newcommand{\N}{\mathbb{N}}
\newcommand{\Z}{\mathbb{Z}}
\newcommand{\Sph}{\mathbb{S}}
\newcommand{\Tor}{\mathbb{T}}
\newcommand{\Id}{\mathrm{Id}}
\newcommand{\Tr}{\operatorname{Tr}}
\newcommand{\rank}{\operatorname{rank}}
\newcommand{\supp}{\operatorname{supp}}
\newcommand{\Lip}{\operatorname{Lip}}
\newcommand{\BV}{\operatorname{BV}}
\newcommand{\TV}{\operatorname{TV}}
\newcommand{\Per}{\operatorname{Per}}
\newcommand{\cutnorm}[1]{\|#1\|_{\square}}
\newcommand{\norm}[1]{\left\|#1\right\|}
\newcommand{\inner}[2]{\left\langle #1, #2 \right\rangle}
\newcommand{\dif}{\,\mathrm{d}}
\newcommand{\abs}[1]{\left|#1\right|}
\newcommand{\WF}{\operatorname{WF}}
\newcommand{\QFI}{\mathrm{QFI}}
\newcommand{\Ric}{\operatorname{Ric}}
\newcommand{\Hess}{\operatorname{Hess}}
\newcommand{\diam}{\operatorname{diam}}
\newcommand{\dist}{\operatorname{dist}}

\begin{document}

\title[Beyond the Spectrum II]{Beyond the Spectrum II: Metric Measure Geometry, Persistent Homology, and Non-Commutative Invariants of Functional Tensor Manifolds}

\author{Reinaldo M. Silva-Filho}
\address{Master's student, Postgraduate Program in Statistics and Agricultural Experimentation (PPGEE/DES), Department of Statistics (DES), Federal University of Lavras (UFLA), Lavras, MG, Brazil}
\email{reinaldo.filho1@estudante.ufla.br}
\date{October 6, 2026}

\begin{abstract}
A functional realization $\Phi$ turns a matrix $A$ into a function $\Phi(A)$ on a domain $\Omega$. Invariants of $\Phi(A)$ can depend on the positions of the entries, which eigenvalues cannot. We study six families of such invariants and mark which results are proved here and which are classical.
(1) For probability-density realizations, the quadratic Wasserstein distance compares matrices of different sizes. We prove the sharp bound $\mathcal{W}_2 \le 2^{-1/2}\diam(\Omega)\,\|\Phi(A)-\Phi(B)\|_{L^1}^{1/2}$ and, for the kernel realization with density bounded below by $c_0$, the metric-speed bound $|\dot\mu| \le C c_0^{-1/2}\|\dot A\|_F$. Bakry--\'Emery curvature gives the classical displacement-convexity and log-Sobolev consequences.
(2) For step realizations we compute super-level persistence of explicit isospectral pairs: a $3\times 3$ pair with different diagrams and equal persistent entropy, and a $5\times 5$ pair whose entropies differ. Persistent entropy depends only on bar lengths and does not separate isospectral matrices in general.
(3) For a trigonometric checkerboard profile, the integrated Willmore energy of level sets equals $c_*\,\epsilon k^3(1+O(1/k))$.
(4) A jump across a smooth interface has wavefront set equal to the conormal bundle (classical); non-constant step realizations have Besov index $s^*_p = 1/p$ for $p > 1$.
(5) On the flat torus, Connes' trace theorem gives $\Tr_\omega(\Phi(A)|\mathcal{D}|^{-d}) = c_d\int\Phi(A)$. The cocycle $\tau_2$ without grading is not topological; with the grading, its antisymmetrization equals $-4i\deg F$.
(6) The quantum Fisher information volume of a field of density matrices vanishes if and only if its partial derivatives are everywhere linearly dependent.
\end{abstract}

\keywords{Functional realizations of matrices, optimal transport, Wasserstein distance, Bakry--\'Emery curvature, persistent homology, persistent entropy, Willmore energy, wavefront set, Besov regularity, Dixmier trace, cyclic cocycles, quantum Fisher information.}
\subjclass[2020]{15A69, 49Q22, 55N31, 53C21, 58J40, 58B34, 81P45, 46E35}

\maketitle

\tableofcontents

\newpage

% =========================================================================
\section{Introduction}
\label{sec:intro}
% =========================================================================

Algebraic invariants of a matrix $A \in \R^{n \times n}$---trace, determinant, spectrum $\sigma(A)$, singular values, rank---are unchanged under simultaneous permutation of rows and columns: for a permutation matrix $P$,
\begin{equation}
    \sigma(PAP^T) = \sigma(A), \quad \Tr(PAP^T) = \Tr(A), \quad \|PAP^T\|_F = \|A\|_F.
\end{equation}
When the indices carry a spatial meaning (pixels of an image, sites of a lattice), these invariants cannot see where an entry sits. That a spectrum does not determine the underlying object is classical: Kac asked whether the Dirichlet spectrum of a planar domain determines its shape \cite{kac1966can}, Gordon, Webb and Wolpert constructed isospectral non-isometric planar domains \cite{gordon1992one}, and van Dam and Haemers survey which graphs are, and which are not, determined by their spectra \cite{vandam2003which}. A \emph{functional realization} is a map
\begin{equation}
    \Phi \colon \R^{n \times n} \supseteq \mathcal{D}_n \longrightarrow \mathcal{F}(\Omega)
\end{equation}
from a class of matrices into a space of functions on a domain $\Omega \subset \R^d$; invariants of the function $\Phi(A)$ may depend on the positions of entries. Volume~I \cite{silvafilho2026beyond} of this series studies first-order invariants such as the Dirichlet energy and the total variation of $\Phi(A)$; functional realizations of matrices and tensors are also the subject of Part~I of the monograph \cite{silvafilho2026geometry}. This volume studies six further families: optimal transport and Bakry--\'Emery curvature (Section~\ref{sec:optimal_transport}), persistent homology (Section~\ref{sec:persistent_homology}), level-set curvature energies and isoperimetry (Section~\ref{sec:willmore_gmt}), wavefront sets and Besov regularity (Section~\ref{sec:microlocal}), Dixmier traces and cyclic cocycles (Section~\ref{sec:noncommutative}), and the quantum Fisher information of fields of density matrices (Section~\ref{sec:quantum_info}). Section~\ref{sec:synthesis} summarizes the results.

Every invariant below depends on the choice of realization. In particular, the Wasserstein distance of Section~\ref{sec:optimal_transport} compares matrices of different sizes only through a chosen realization for each size; other ways of comparing objects of different sizes, such as the cut distance of graph limits \cite{lovasz2012large}, Sturm's transportation distance between metric measure spaces \cite{sturm2006geometryI}, the Gromov--Wasserstein distance \cite{memoli2011gromov}, and its extension by Chowdhury and M\'emoli to weighted networks with different numbers of nodes \cite{chowdhury2019gromov}, are not discussed here. Sturm showed that the space of all metric measure spaces with the $L^{2,q}$-distortion distance has non-negative curvature in the sense of Alexandrov \cite{sturm2023space}.

\subsection*{Related work}
Each section draws on an established theory. Optimal transport and its dynamic formulation are treated in \cite{villani2009optimal,ambrosio2008gradient,benamou2000computational}; for centred Gaussian measures with covariance matrices $A$ and $B$ the quadratic Wasserstein distance is the Bures--Wasserstein distance between $A$ and $B$ \cite{bhatia2019bures}, a transport distance on positive definite matrices of a fixed size. Lower Ricci bounds defined through displacement convexity of the entropy go back to Sturm and to Lott and Villani \cite{sturm2006geometryI,sturm2006geometryII,lott2009ricci}, and the Bakry--\'Emery calculus is developed in \cite{bakry1985diffusions,bakry2014analysis}. Persistent homology was introduced in \cite{edelsbrunner2002topological,zomorodian2005computing}; see \cite{edelsbrunner2010computational} for a textbook account, \cite{kaczynski2004computational,wagner2012efficient} for cubical complexes such as the super-level sets of step realizations, and \cite{chintakunta2015entropy,atienza2020stability} for persistent entropy. The level-set curvature energy of Section~\ref{sec:willmore_gmt} is, up to the weights, the curvature term of the elastica-type energies used in image inpainting \cite{chan2003euler,ambrosio2003direct}. For wavefront sets we follow H\"ormander \cite{hormander2003analysis}; for the non-commutative part, Connes \cite{connes1985noncommutative,connes1988action,connes1994noncommutative}, Connes and Moscovici \cite{connes1995local}, and the monograph \cite{lord2013singular} on singular traces. The quantum Fisher information goes back to Helstrom \cite{helstrom1967minimum} and Braunstein and Caves \cite{braunstein1994statistical}; see \cite{liu2020quantum} for a review.

\subsection*{What is new here}
Cited classical results are marked ``classical'' in their headings and are applied after checking their hypotheses; remarks that rest only on numerical evidence say so. Our own contributions are mostly explicit computations: the bounds of Theorem~\ref{thm:wasserstein_properties}, the isospectral examples of Theorem~\ref{thm:isospectral_separation}, the scaling law of Proposition~\ref{prop:willmore_scaling}, the Besov indices of Proposition~\ref{prop:besov}, the cocycle formulas of Proposition~\ref{prop:cocycle}, and the rank criterion of Theorem~\ref{thm:qfi_volume}.

% =========================================================================
\section{Optimal Transport and Bakry--\'Emery Curvature}
\label{sec:optimal_transport}
% =========================================================================

\subsection{The Density Realization Operator}
Let $V$ be an $n$-dimensional Euclidean space. Let $\mathcal{S}_+^n \subset \R^{n \times n}$ denote the cone of symmetric positive semi-definite (PSD) matrices, normalized such that $\Tr(A) = 1$. Let $\Omega = [0, 1]^d$ be the unit hypercube equipped with the standard Lebesgue measure $\dif x$.

\begin{definition}[Probability Density Realization]
Let $\Phi \colon \mathcal{S}_+^n \to L^1(\Omega)$ be a positive functional realization operator satisfying:
\begin{equation}
    \Phi(A)(x) \ge 0 \quad \text{a.e. } x \in \Omega, \quad \text{and} \quad \int_\Omega \Phi(A)(x) \dif x = \Tr(A) = 1.
\end{equation}
Each normalized matrix $A \in \mathcal{S}_+^n$ thus induces an absolutely continuous Borel probability measure on $\Omega$:
\begin{equation}
    \dif \mu_A(x) \coloneqq \Phi(A)(x) \dif x \in \mathcal{P}_{\mathrm{ac}}(\Omega).
\end{equation}
\end{definition}

A family of such realizations is the \emph{kernel realization}. Fix real functions $\psi_1, \dots, \psi_n \in L^4(\Omega)$ that are orthonormal in $L^2(\Omega)$, $\int_\Omega \psi_i \psi_j \dif x = \delta_{ij}$, and set
\begin{equation}
    \Phi_{\mathrm{kernel}}(A)(x) \coloneqq \sum_{i,j=1}^n A_{ij} \psi_i(x) \psi_j(x) = \boldsymbol{\psi}(x)^T A \boldsymbol{\psi}(x).
    \label{eq:kernel_realization}
\end{equation}
Since $A \succeq 0$, $\Phi_{\mathrm{kernel}}(A) \ge 0$, and orthonormality gives $\int_\Omega \Phi_{\mathrm{kernel}}(A) \dif x = \sum_i A_{ii} = \Tr(A) = 1$. The map $A \mapsto \Phi_{\mathrm{kernel}}(A)$ is linear, and since $|\boldsymbol{\psi}^T M \boldsymbol{\psi}| \le \|M\|_{\mathrm{op}} |\boldsymbol{\psi}|^2 \le \|M\|_F |\boldsymbol{\psi}|^2$ for symmetric $M$,
\begin{equation}
    \|\Phi_{\mathrm{kernel}}(M)\|_{L^2(\Omega)} \le L_\psi \|M\|_F, \qquad L_\psi \coloneqq \Big\| \sum_{i=1}^n \psi_i^2 \Big\|_{L^2(\Omega)}.
    \label{eq:kernel_lipschitz}
\end{equation}

\subsection{Wasserstein Distance Between Matrices of Different Sizes}
Let $A \in \mathcal{S}_+^n$ and $B \in \mathcal{S}_+^m$, possibly with $n \ne m$, and fix a density realization for each size. The realizations induce probability measures $\mu_A, \mu_B \in \mathcal{P}(\Omega)$.

\begin{definition}[Functional $L^2$-Wasserstein Matrix Distance]
The functional quadratic Wasserstein distance between $A \in \mathcal{S}_+^n$ and $B \in \mathcal{S}_+^m$ is defined by:
\begin{equation}
    \mathcal{W}_2(A, B) \coloneqq \left( \inf_{\pi \in \Pi(\mu_A, \mu_B)} \int_{\Omega \times \Omega} \|x - y\|^2 \dif \pi(x, y) \right)^{1/2},
    \label{eq:wasserstein_matrix}
\end{equation}
where $\Pi(\mu_A, \mu_B)$ is the set of all Borel probability measures on $\Omega \times \Omega$ having marginals $\mu_A$ and $\mu_B$.
\end{definition}
This is Kantorovich's relaxation \cite{kantorovich2006translocation} of the transport problem, with quadratic cost; see \cite[Ch.~6]{villani2009optimal}.

By the Benamou--Brenier formula \cite{benamou2000computational}, $\mathcal{W}_2(A, B)$ also has a dynamic formulation:
\begin{equation}
    \mathcal{W}_2^2(A, B) = \inf_{(\rho_t, \mathbf{v}_t)} \left\{ \int_0^1 \int_\Omega \|\mathbf{v}_t(x)\|^2 \rho_t(x) \dif x \dif t \right\},
\end{equation}
subject to the continuity equation and boundary conditions:
\begin{equation}
    \partial_t \rho_t + \nabla \cdot (\rho_t \mathbf{v}_t) = 0, \quad \rho_0 = \Phi(A), \quad \rho_1 = \Phi(B).
\end{equation}

\begin{theorem}[Metric properties]
\label{thm:wasserstein_properties}
Let $\mathcal{M} = \bigcup_{n=1}^\infty \mathcal{S}_+^n$, with a density realization fixed for each $n$, and $\Omega = [0,1]^d$. Then:
\begin{enumerate}[label=(\alph*)]
    \item $\mathcal{W}_2$ is a pseudometric on $\mathcal{M}$, and $\mathcal{W}_2(A,B) = 0$ if and only if $\Phi(A) = \Phi(B)$ almost everywhere. Hence $\mathcal{W}_2$ is a metric on the quotient of $\mathcal{M}$ by the relation $\Phi(A) = \Phi(B)$ a.e.
    \item For all $A, B \in \mathcal{M}$,
    \begin{equation}
        \mathcal{W}_2(A, B) \le \frac{\diam(\Omega)}{\sqrt{2}} \|\Phi(A) - \Phi(B)\|_{L^1(\Omega)}^{1/2},
        \label{eq:w2_l1}
    \end{equation}
    and the constant $\diam(\Omega)/\sqrt{2}$ cannot be decreased.
    \item Let $\Phi = \Phi_{\mathrm{kernel}}$ be the kernel realization \eqref{eq:kernel_realization}, let $\tau \mapsto A(\tau) \in \mathcal{S}_+^n$ be a $C^1$ curve on an open interval $I$, and assume $\rho_\tau \coloneqq \Phi(A(\tau)) \ge c_0 > 0$ on $\Omega$ for all $\tau \in I$. Then $\tau \mapsto \mu_{A(\tau)}$ is absolutely continuous in $(\mathcal{P}_2(\Omega), \mathcal{W}_2)$, and for almost every $\tau$ its metric derivative satisfies
    \begin{equation}
        |\dot{\mu}|(\tau) \coloneqq \lim_{h \to 0} \frac{\mathcal{W}_2(A(\tau+h), A(\tau))}{|h|} \le \frac{\|\partial_\tau \rho_\tau\|_{L^2(\Omega)}}{\pi \sqrt{c_0}} \le \frac{L_\psi}{\pi \sqrt{c_0}} \|\dot{A}(\tau)\|_F,
        \label{eq:metric_speed}
    \end{equation}
    with $L_\psi$ as in \eqref{eq:kernel_lipschitz}.
\end{enumerate}
\end{theorem}

\begin{proof}
(a) $\mathcal{W}_2$ is a metric on $\mathcal{P}(\Omega)$ for compact $\Omega$ \cite[Ch.~6]{villani2009optimal}. Its pull-back by $A \mapsto \mu_A$ is therefore a pseudometric that vanishes exactly when $\mu_A = \mu_B$, that is, when $\Phi(A) = \Phi(B)$ a.e.

(b) Write $\rho = \Phi(A)$, $\sigma = \Phi(B)$, $m \coloneqq \frac{1}{2}\|\rho - \sigma\|_{L^1} = \int (\rho - \sigma)_+ = \int (\sigma - \rho)_+$. If $m = 0$ there is nothing to prove. Otherwise the measure
\begin{equation}
    \pi \coloneqq (\mathrm{id}, \mathrm{id})_\# \big(\min(\rho, \sigma) \dif x\big) + \frac{1}{m} (\rho - \sigma)_+ \dif x \otimes (\sigma - \rho)_+ \dif y
\end{equation}
has marginals $\mu_A$ and $\mu_B$. The first part costs nothing and the second has total mass $m$, so $\mathcal{W}_2^2(A,B) \le \int \|x-y\|^2 \dif \pi \le \diam(\Omega)^2 m$, which is \eqref{eq:w2_l1}. For sharpness, take $d = 1$, the kernel realization with $\psi_1 = \epsilon^{-1/2}\mathbf{1}_{[0,\epsilon]}$, $\psi_2 = \epsilon^{-1/2}\mathbf{1}_{[1-\epsilon,1]}$ ($0 < \epsilon < 1/2$), and $A = \mathrm{diag}(1,0)$, $B = \mathrm{diag}(0,1)$, so that $\Phi(A) = \epsilon^{-1}\mathbf{1}_{[0,\epsilon]}$ and $\Phi(B) = \epsilon^{-1}\mathbf{1}_{[1-\epsilon,1]}$. The monotone map $x \mapsto x + 1 - \epsilon$ is optimal, so $\mathcal{W}_2 = 1 - \epsilon$, while the right-hand side of \eqref{eq:w2_l1} equals $1$. In dimension $d$, the same construction near two opposite corners gives ratios tending to $1$.

(c) Since $\Phi$ is linear, $\partial_\tau \rho_\tau = \Phi(\dot{A}(\tau))$ exists in $L^2(\Omega)$, depends continuously on $\tau$, and $\|\partial_\tau \rho_\tau\|_{L^2} \le L_\psi \|\dot{A}(\tau)\|_F$ by \eqref{eq:kernel_lipschitz}. Because $\Tr A(\tau) = 1$, $\int_\Omega \partial_\tau \rho_\tau = 0$. By the Lax--Milgram lemma there is a unique $\varphi_\tau \in H^1(\Omega)$ with $\int_\Omega \varphi_\tau = 0$ and
\begin{equation}
    \int_\Omega \nabla \varphi_\tau \cdot \nabla \zeta \dif x = \int_\Omega \partial_\tau \rho_\tau \, \zeta \dif x \qquad \text{for all } \zeta \in H^1(\Omega),
    \label{eq:weak_neumann}
\end{equation}
the weak form of $-\Delta \varphi_\tau = \partial_\tau\rho_\tau$ with Neumann data. The first non-zero Neumann eigenvalue of $-\Delta$ on $[0,1]^d$ is $\pi^2$, so the Poincar\'e inequality $\|\varphi_\tau\|_{L^2} \le \pi^{-1} \|\nabla \varphi_\tau\|_{L^2}$ holds; taking $\zeta = \varphi_\tau$ in \eqref{eq:weak_neumann} gives $\|\nabla \varphi_\tau\|_{L^2} \le \pi^{-1}\|\partial_\tau \rho_\tau\|_{L^2}$. Put $\mathbf{v}_\tau \coloneqq \nabla \varphi_\tau / \rho_\tau$. For every $\zeta \in C_c^\infty(\R^d)$, \eqref{eq:weak_neumann} gives
\begin{equation}
    \frac{\dif}{\dif \tau} \int \zeta \dif \mu_{A(\tau)} = \int_\Omega \zeta \, \partial_\tau \rho_\tau \dif x = \int_\Omega \nabla \zeta \cdot \mathbf{v}_\tau \, \rho_\tau \dif x,
\end{equation}
so $(\mu_{A(\tau)}, \mathbf{v}_\tau)$ solves the continuity equation in the distributional sense on $I \times \R^d$, and
\begin{equation}
    \|\mathbf{v}_\tau\|_{L^2(\mu_{A(\tau)})}^2 = \int_\Omega \frac{|\nabla \varphi_\tau|^2}{\rho_\tau} \dif x \le \frac{1}{c_0} \|\nabla \varphi_\tau\|_{L^2}^2 \le \frac{\|\partial_\tau \rho_\tau\|_{L^2}^2}{\pi^2 c_0}.
\end{equation}
This bound is locally bounded in $\tau$. By \cite[Thm.~8.3.1]{ambrosio2008gradient}, a narrowly continuous curve solving the continuity equation with velocity field $\mathbf{v}_\tau$ satisfying $\int_I \|\mathbf{v}_\tau\|_{L^2(\mu_\tau)} \dif\tau < \infty$ on compact subintervals is absolutely continuous, and $|\dot{\mu}|(\tau) \le \|\mathbf{v}_\tau\|_{L^2(\mu_\tau)}$ for a.e.\ $\tau$. This gives \eqref{eq:metric_speed}.
\end{proof}

\begin{remark}
The constant $1/\pi$ in the first inequality of \eqref{eq:metric_speed} is attained in the limit: for $d = 1$ and $\rho_\tau = 1 + \tau \epsilon \cos(\pi x)$ at $\tau = 0$, the velocity field above is the exact one-dimensional transport velocity, and the ratio of the two sides tends to $1$. A numerical check with exact one-dimensional Wasserstein distances confirms \eqref{eq:metric_speed} on several curves. The first inequality in \eqref{eq:metric_speed} is an instance of the formal identity between $\mathcal{W}_2$ and a weighted homogeneous $\dot H^{-1}$ norm for infinitesimal perturbations; Peyre \cite{peyre2018comparison} integrates this identity into the non-asymptotic bound $\mathcal{W}_2(\mu,\nu) \le 2\|\mu - \nu\|_{\dot H^{-1}(\mu)}$.
\end{remark}

\subsection{Bakry--\'Emery Metric Curvature of a Matrix Realization}
Given a strictly positive functional realization $\Phi(A) \in C^2(\Omega)$ with $\Phi(A)(x) > 0$, we view $(\Omega, g, \dif \mu_A)$ as a weighted Riemannian metric measure space (smooth metric measure space) with background Euclidean metric $g_{ij} = \delta_{ij}$ and smooth potential:
\begin{equation}
    V_A(x) \coloneqq -\log \Phi(A)(x) \implies \dif \mu_A(x) = e^{-V_A(x)} \dif x.
\end{equation}
The weighted Laplacian (drift Laplacian / Witten Laplacian) acting on $f \in C^2(\Omega)$ is:
\begin{equation}
    \Delta_A f \coloneqq \Delta f - \langle \nabla V_A, \nabla f \rangle = \Delta f + \frac{\langle \nabla \Phi(A), \nabla f \rangle}{\Phi(A)}.
\end{equation}

\begin{definition}[Bakry--\'Emery Curvature Tensor and Invariant]
The infinite-dimensional Bakry--\'Emery Ricci tensor associated with the matrix $A$ is:
\begin{equation}
    \Ric_\infty(\Phi(A)) \coloneqq \Ric_g + \nabla^2 V_A = -\nabla^2 \log \Phi(A) = -\frac{\nabla^2 \Phi(A)}{\Phi(A)} + \frac{\nabla \Phi(A) \otimes \nabla \Phi(A)}{\Phi(A)^2}.
\end{equation}
The \emph{Bakry--\'Emery Curvature Invariant} $\kappa_{\mathrm{BE}}(A)$ of the matrix is the infimum of its spectrum:
\begin{equation}
    \kappa_{\mathrm{BE}}(A) \coloneqq \inf_{x \in \Omega, \|\xi\|=1} \left\langle \left( -\nabla^2 \log \Phi(A)(x) \right) \xi, \xi \right\rangle.
\end{equation}
\end{definition}

\begin{theorem}[Bakry--\'Emery consequences; (b) and (c) classical]
\label{thm:bakry_emery}
Let $\Omega = [0, 1]^d$ and let $A \in \mathcal{S}_+^n$ have a realization $\Phi(A) \in C^2(\overline{\Omega})$ with $\Phi(A) > 0$ on $\overline{\Omega}$ and $\kappa_{\mathrm{BE}}(A) \ge K > 0$. Let $\Delta_A$ carry Neumann boundary conditions, that is, let $e^{\tau\Delta_A}$ be the self-adjoint semigroup on $L^2(\mu_A)$ associated with the Dirichlet form $f \mapsto \int_\Omega \|\nabla f\|^2 \dif\mu_A$ on $H^1(\Omega)$. Then:
\begin{enumerate}[label=(\alph*)]
    \item \textbf{Displacement convexity.} For $\mu_0, \mu_1 \in \mathcal{P}(\Omega)$ absolutely continuous with finite entropy, let $(\mu_t)_{t \in [0,1]}$ be their displacement interpolation. Then $\mathrm{Ent}(\mu \mid \mu_A) \coloneqq \int_\Omega \rho \log(\rho / \Phi(A)) \dif x$ satisfies
    \begin{equation}
        \mathrm{Ent}(\mu_t \mid \mu_A) \le (1-t)\,\mathrm{Ent}(\mu_0 \mid \mu_A) + t\,\mathrm{Ent}(\mu_1 \mid \mu_A) - \frac{K}{2} t(1-t) \mathcal{W}_2^2(\mu_0, \mu_1).
    \end{equation}
    \item \textbf{Logarithmic Sobolev inequality.} For every $f \in H^1(\Omega)$,
    \begin{equation}
        \int_\Omega f^2 \log\left(\frac{f^2}{\int_\Omega f^2 \dif\mu_A}\right) \dif\mu_A \le \frac{2}{K} \int_\Omega \|\nabla f\|^2 \dif\mu_A.
    \end{equation}
    \item \textbf{Variance decay.} For every $u_0 \in L^2(\mu_A)$, $e^{\tau\Delta_A}u_0$ converges to the constant $\int u_0 \dif\mu_A$ and
    \begin{equation}
        \mathrm{Var}_{\mu_A}(e^{\tau \Delta_A} u_0) \le e^{-2K\tau} \mathrm{Var}_{\mu_A}(u_0).
    \end{equation}
\end{enumerate}
\end{theorem}

\begin{proof}
(a) Write $V = V_A = -\log \Phi(A)$, so $\nabla^2 V \ge K\,\Id$ on the convex set $\Omega$, and $\mathrm{Ent}(\mu \mid \mu_A) = \int \rho \log \rho \dif x + \int V \dif \mu$. By Brenier's theorem \cite{brenier1991polar} there is an optimal map $T$ with $\mu_1 = T_\#\mu_0$, and $\mu_t = ((1-t)\mathrm{id} + tT)_\# \mu_0$; since $\Omega$ is convex, $\mu_t$ is supported in $\Omega$. McCann's displacement convexity theorem \cite{mccann1997convexity} gives $\int \rho_t \log \rho_t \le (1-t)\int \rho_0\log\rho_0 + t \int\rho_1\log\rho_1$. For the potential term, $K$-convexity of $V$ along the segment from $x$ to $T(x)$ gives
\begin{equation}
    V((1-t)x + tT(x)) \le (1-t)V(x) + tV(T(x)) - \frac{K}{2}t(1-t)\|x - T(x)\|^2,
\end{equation}
and integrating against $\mu_0$, with $\int \|x - T(x)\|^2 \dif\mu_0 = \mathcal{W}_2^2(\mu_0,\mu_1)$, gives the claim.

(b), (c) These are classical consequences of the Bakry--\'Emery criterion \cite{bakry1985diffusions}; see \cite{bakry2014analysis} for the theory. On the convex body $\Omega$ the measure $\mu_A$ has the form $f\,\dif\gamma$, where $\gamma$ is the Gaussian measure with covariance $K^{-1}\Id$ and $f = Z^{-1}\mathbf{1}_\Omega\, e^{-(V - K|x|^2/2)}$ is log-concave. By Caffarelli's contraction theorem \cite{caffarelli2000monotonicity}, the Brenier map from $\gamma$ to $\mu_A$ is $1$-Lipschitz, so $\mu_A$ inherits the Gaussian logarithmic Sobolev inequality with constant $2/K$, which is (b). Linearizing (b) gives the Poincar\'e inequality $\mathrm{Var}_{\mu_A}(f) \le K^{-1}\int\|\nabla f\|^2\dif\mu_A$, so the spectral gap of $-\Delta_A$ is at least $K$, and (c) follows from the spectral theorem. We do not reprove these classical statements.
\end{proof}

\begin{remark}
On smooth Riemannian manifolds without boundary and without weight, lower Ricci bounds are equivalent to displacement convexity of the entropy \cite{von2005transport}. Sturm \cite{sturm2006geometryI,sturm2006geometryII} and Lott and Villani \cite{lott2009ricci} turned this property into a definition of lower Ricci bounds for metric measure spaces. Related links are the identification of the Fokker--Planck equation as the Wasserstein gradient flow of the relative entropy \cite{jordan1998variational} and the implication from a logarithmic Sobolev inequality to Talagrand's transport inequality \cite{otto2000generalization}. None of these results is used above.
\end{remark}

\begin{example}[Realizations with Strict Positive Curvature]
\label{ex:positive_curvature}
For $n = d$, a family of realizations with $\kappa_{\mathrm{BE}}(A) > 0$ is the truncated Gaussian realization on $\Omega = [0, 1]^d$, with a fixed centre $x_0$ and $Z_A$ the normalizing constant:
\begin{equation}
    \Phi_{\mathrm{Gibbs}}(A)(x) \coloneqq \frac{1}{Z_A} \exp\left( -\frac{1}{2} (x - x_0)^T A (x - x_0) \right), \quad A \succ 0.
\end{equation}
The potential is $V_A(x) = \frac{1}{2} (x - x_0)^T A (x - x_0) + \log Z_A$, yielding $\nabla^2 V_A(x) = A$. Therefore, $\Ric_\infty(\Phi_{\mathrm{Gibbs}}(A)) = A \ge \lambda_{\min}(A) \cdot \Id$, and the Bakry--\'Emery curvature invariant is precisely the smallest eigenvalue of the matrix:
\begin{equation}
    \kappa_{\mathrm{BE}}(A) = \lambda_{\min}(A) > 0.
\end{equation}
Thus every positive definite $A$ has $\kappa_{\mathrm{BE}}(A) > 0$ under this realization. Under the normalization $\Tr A = 1$ one has $\kappa_{\mathrm{BE}}(A) = \lambda_{\min}(A) \le 1/d$.
\end{example}

% =========================================================================
\section{Persistent Homology and Persistent Entropy}
\label{sec:persistent_homology}
% =========================================================================

\subsection{The Super-Level Set Filtration}
Let $A \in \R^{n \times n}$ and let $\Phi(A) \colon \Omega \to \R$ be a realization on $\Omega = [0, 1]^d$ that is either continuous or a step realization (Definition~\ref{def:step}).

\begin{definition}[Super-level filtration]
For $t \in \R$, the super-level set of $A$ is
\begin{equation}
    X_t(A) \coloneqq \left\{ x \in \Omega \;\middle|\; \Phi(A)(x) \ge t \right\}.
\end{equation}
For $s \ge t$ the inclusion $X_s(A) \hookrightarrow X_t(A)$ induces maps in homology, giving a persistence module indexed by $(\R, \ge)$:
\begin{equation}
    \mathcal{H}_k(A) \colon t \longmapsto H_k(X_t(A); \mathbb{F}),
\end{equation}
where $\mathbb{F}$ is a field and $H_k$ is singular homology.
\end{definition}
Persistence of a filtration was introduced by Edelsbrunner, Letscher and Zomorodian \cite{edelsbrunner2002topological} and given its module-theoretic form by Zomorodian and Carlsson \cite{zomorodian2005computing}; see also \cite{edelsbrunner2010computational}.

\begin{definition}[Step realization]
\label{def:step}
Let $d = 2$, $A \in \R^{n \times n}$, and let $\Omega_{ij} = [\frac{i-1}{n}, \frac{i}{n}] \times [\frac{j-1}{n}, \frac{j}{n}]$ be the closed cells of the uniform $n \times n$ grid. The step realization is
\begin{equation}
    \Phi_\square(A)(x) \coloneqq \max\{ A_{ij} : x \in \Omega_{ij} \}.
\end{equation}
It is upper semicontinuous, and $X_t(A) = \bigcup\{\Omega_{ij} : A_{ij} \ge t\}$ is a finite union of closed cells.
\end{definition}
The super-level sets of a step realization are cubical complexes, as for grey-scale images; for the homology and persistence of cubical complexes see \cite{kaczynski2004computational,wagner2012efficient}.

When every $H_k(X_t(A))$ is finite-dimensional---for instance for step realizations, where $X_t(A)$ is a finite cubical complex taking finitely many values---the module decomposes into interval modules \cite{crawley2015decomposition}. A bar on which the class is alive for $t \in (d_i, b_i]$ is recorded as the point $(b_i, d_i)$, $b_i > d_i$: $b_i$ is the threshold at which the class appears and $d_i$ the threshold at which it dies as $t$ decreases. The \emph{persistence diagram} $\mathrm{Dgm}_k(\Phi(A))$ is the multiset of these points together with the diagonal.

\subsection{Bottleneck Stability of Functional Tensors}
Let $\mathrm{Dgm}_k(f)$ and $\mathrm{Dgm}_k(g)$ be two persistence diagrams. The bottleneck distance is:
\begin{equation}
    d_B(\mathrm{Dgm}_k(f), \mathrm{Dgm}_k(g)) \coloneqq \inf_{\gamma \colon \mathrm{Dgm}_k(f) \to \mathrm{Dgm}_k(g)} \sup_{p \in \mathrm{Dgm}_k(f)} \|p - \gamma(p)\|_\infty,
\end{equation}
where $\gamma$ ranges over all bijections between the multisets.

\begin{theorem}[Bottleneck stability; classical]
\label{thm:bottleneck_stability}
Let $A, B \in \R^{n \times n}$ and let $\Phi$ be a realization operator with values in $C(\Omega)$ and Lipschitz constant $L_\Phi$:
\begin{equation}
    \|\Phi(A) - \Phi(B)\|_{L^\infty(\Omega)} \le L_\Phi \|A - B\|_F.
\end{equation}
Assume $\Phi(A)$ and $\Phi(B)$ are tame continuous functions (e.g., Morse functions with finitely many critical values, or piecewise affine functions on a finite simplicial triangulation of $\Omega$), ensuring well-defined, finite persistence diagrams. Then, for every homology dimension $k \in \{0, 1, \dots, d-1\}$, the persistence diagrams satisfy:
\begin{equation}
    d_B(\mathrm{Dgm}_k(\Phi(A)), \mathrm{Dgm}_k(\Phi(B))) \le \|\Phi(A) - \Phi(B)\|_{L^\infty(\Omega)} \le L_\Phi \|A - B\|_F.
\end{equation}
\end{theorem}

\begin{proof}
Let $f = \Phi(A)$ and $g = \Phi(B)$. For any $\epsilon = \|f - g\|_{L^\infty(\Omega)}$, we have $f(x) \le g(x) + \epsilon$ and $g(x) \le f(x) + \epsilon$ for all $x \in \Omega$. 
This induces interlacing inclusions between the super-level sets:
\begin{equation}
    X_{t+\epsilon}(f) \subseteq X_t(g) \subseteq X_{t-\epsilon}(f).
\end{equation}
Applying $H_k(\cdot; \mathbb{F})$ produces an $\epsilon$-interleaving of the persistence modules $\mathcal{H}_k(f)$ and $\mathcal{H}_k(g)$. By the algebraic stability theorem \cite{chazal2009proximity,chazal2016structure}, an $\epsilon$-interleaving of such modules bounds the bottleneck distance of their diagrams by $\epsilon$:
\begin{equation}
    d_B(\mathrm{Dgm}_k(f), \mathrm{Dgm}_k(g)) \le \epsilon = \|f - g\|_{L^\infty(\Omega)}.
\end{equation}
For tame functions this is the stability theorem of Cohen-Steiner, Edelsbrunner and Harer \cite{cohensteiner2007stability}. Combining it with the Lipschitz bound on $\Phi$ completes the proof.
\end{proof}

\begin{remark}
Theorem~\ref{thm:bottleneck_stability} needs continuous realizations and does not apply to the step realizations of Definition~\ref{def:step}. Nor can a step realization be approximated uniformly by mollified ones: a uniform limit of continuous functions is continuous, so $\|\Phi_\square(A) * \eta_\epsilon - \Phi_\square(A)\|_{L^\infty} \not\to 0$ whenever $\Phi_\square(A)$ has a jump.
\end{remark}

\subsection{The Persistent Entropy Invariant}
\begin{definition}[Persistent Entropy of a Matrix]
Let $\mathrm{Dgm}_k(\Phi(A)) = \{(b_i, d_i)\}_{i=1}^{N_k}$ be the off-diagonal points of the $k$-th persistence diagram. The lifespan of the $i$-th topological feature is $\ell_i \coloneqq b_i - d_i > 0$. The total persistent lifetime is $L_k(A) \coloneqq \sum_{i=1}^{N_k} \ell_i$. The persistent probability weights are $p_i \coloneqq \ell_i / L_k(A)$.
The \emph{$k$-th Persistent Entropy} of the matrix is:
\begin{equation}
    E_{\mathrm{pers}}^{(k)}(A) \coloneqq -\sum_{i=1}^{N_k} p_i \log p_i \ge 0.
\end{equation}
By convention, if $N_k = 0$ (i.e., the persistence diagram $\mathrm{Dgm}_k(\Phi(A))$ is empty and has no off-diagonal bars), we set $E_{\mathrm{pers}}^{(k)}(A) \coloneqq 0$.
\end{definition}
Persistent entropy, the Shannon entropy of the normalized bar lengths, was introduced by Chintakunta, Gentimis, Gonzalez-D{\'\i}az, Jim{\'e}nez and Krim \cite{chintakunta2015entropy}; Atienza, Gonzalez-D{\'\i}az and Soriano-Trigueros proved that it is stable under small perturbations of the barcode \cite{atienza2020stability}.

\begin{theorem}[Isospectral pairs separated by step persistence]
\label{thm:isospectral_separation}
Use the step realization $\Phi_\square$ of Definition~\ref{def:step} and homology with coefficients in a field $\mathbb{F}$.
\begin{enumerate}[label=(\alph*)]
    \item Let $B$ be the $3 \times 3$ matrix with $B_{22} = 0$ and all other entries $2$, let $P$ be the permutation matrix exchanging the indices $1$ and $2$, and let $A = PBP^T$. Then $A$ and $B$ are symmetric and isospectral, $\sigma(A) = \sigma(B) = \{2 + 2\sqrt{3}, 2 - 2\sqrt{3}, 0\}$, and
    \begin{equation}
        \mathrm{Dgm}_1(\Phi_\square(A)) = \emptyset, \qquad \mathrm{Dgm}_1(\Phi_\square(B)) = \{(2,0)\}, \qquad E^{(1)}_{\mathrm{pers}}(A) = E^{(1)}_{\mathrm{pers}}(B) = 0.
    \end{equation}
    \item Let $B'$ be the $5 \times 5$ matrix with $B'_{22} = 0$, $B'_{44} = 1$ and all other entries $3$, let $P'$ exchange the indices $1$ and $2$, and let $A' = P'B'P'^T$. Then $A'$ and $B'$ are symmetric and isospectral,
    \begin{equation}
        \mathrm{Dgm}_1(\Phi_\square(B')) = \{(3,0), (3,1)\}, \qquad \mathrm{Dgm}_1(\Phi_\square(A')) = \{(3,1)\},
    \end{equation}
    and $E^{(1)}_{\mathrm{pers}}(B') = -\frac{3}{5}\log\frac{3}{5} - \frac{2}{5}\log\frac{2}{5} \approx 0.673 > 0 = E^{(1)}_{\mathrm{pers}}(A')$.
\end{enumerate}
\end{theorem}

\begin{proof}
We use one homological fact. Let $K \subset \R^2$ be a compact contractible union of closed grid cells and let $U_1, \dots, U_r$ be the interiors of closed grid cells $\overline{U}_1, \dots, \overline{U}_r \subset \operatorname{int} K$ with pairwise disjoint closures. Then $H_1(K \setminus \bigcup_l U_l; \mathbb{F}) \cong \mathbb{F}^r$, with basis the classes of the boundary loops $\partial U_l$. For $r = 1$, choose a closed square $\overline{W}$ with $\overline{U}_1 \subset \operatorname{int}\overline{W}$ and $\overline{W} \subset \operatorname{int} K$, and put $Y \coloneqq K \setminus \operatorname{int}\overline{W}$; all sets involved are subcomplexes of a refined cubical structure. The Mayer--Vietoris sequence of $K = Y \cup \overline{W}$, with $Y \cap \overline{W} = \partial\overline{W}$ a circle and $H_1(K) = H_2(K) = H_1(\overline{W}) = 0$, shows that $Y$ is connected and that $H_1(Y) \cong \mathbb{F}$, generated by $[\partial\overline{W}]$. The Mayer--Vietoris sequence of $K \setminus U_1 = Y \cup (\overline{W} \setminus U_1)$, with intersection $\partial\overline{W}$ and the annulus $\overline{W} \setminus U_1$ having $H_1 \cong \mathbb{F}$ generated by $[\partial\overline{W}] = [\partial U_1]$, then gives $H_1(K \setminus U_1) \cong \mathbb{F}^2 / \mathbb{F}(1,1) \cong \mathbb{F}$, generated by $[\partial U_1]$. The general case follows by induction on $r$, using pairwise disjoint squares $\overline{W}_l$: since $H_2$ of a compact planar cubical complex vanishes, the same two sequences show that removing one more $U_l$ adds one direct summand, generated by $[\partial U_l]$. Naturality of the Mayer--Vietoris sequence shows that adding back the cell $\overline{U}_l$ kills $[\partial U_l]$ and maps the other basis classes to the corresponding basis classes.

(a) $A$ and $B$ are orthogonally similar, hence isospectral; the eigenvalues of $B$ are computed directly. For $t > 2$ both super-level sets are empty, and for $t \le 0$ both are the whole square. For $t \in (0,2]$, $X_t(B)$ is the square minus the open central cell, so $H_1 \cong \mathbb{F}$, and the class dies at $t = 0$: $\mathrm{Dgm}_1(\Phi_\square(B)) = \{(2,0)\}$. Since $A_{11} = 0$ and the other entries equal $2$, $X_t(A)$ for $t \in (0,2]$ is the square minus $[0,\frac13) \times [0, \frac13)$. This set is star-shaped with respect to $(1,1)$: if $p$ lies in it, then $p_1 \ge \frac13$ or $p_2 \ge \frac13$, and that coordinate stays $\ge \frac13$ along the segment from $p$ to $(1,1)$. Hence $H_1(X_t(A)) = 0$ for all $t$ and $\mathrm{Dgm}_1(\Phi_\square(A)) = \emptyset$. A diagram with at most one bar has entropy $0$, so both entropies vanish.

(b) Write $C_{ij}$ for the open cells of the $5 \times 5$ grid. For $B'$: $X_t = \emptyset$ for $t > 3$; for $t \in (1, 3]$, $X_t$ is the square minus $C_{22}$ and $C_{44}$, whose closures are disjoint and lie in the interior, so $H_1 \cong \mathbb{F}^2$ with basis $[\partial C_{22}], [\partial C_{44}]$; for $t \in (0, 1]$, the cell $\overline{C}_{44}$ is added, which kills $[\partial C_{44}]$ and leaves $H_1 \cong \mathbb{F}$ generated by $[\partial C_{22}]$; for $t \le 0$ the square is contractible. The module is therefore the sum of the interval modules on $(0,3]$ and $(1,3]$, so $\mathrm{Dgm}_1 = \{(3,0),(3,1)\}$ with lifetimes $3$ and $2$, weights $\frac35, \frac25$, and the stated entropy. For $A'$, the entries are $A'_{11} = 0$, $A'_{44} = 1$ and all others $3$. For $t \in (1,3]$, $X_t$ is the star-shaped set $K' \coloneqq [0,1]^2 \setminus [0, \frac15)^2$ (same argument as in (a)) minus $C_{44}$, with $\overline{C}_{44} \subset \operatorname{int} K'$, so $H_1 \cong \mathbb{F}$; for $t \in (0,1]$, $X_t = K'$ is contractible. Hence $\mathrm{Dgm}_1(\Phi_\square(A')) = \{(3,1)\}$ and $E^{(1)}_{\mathrm{pers}}(A') = 0$.
\end{proof}

\begin{remark}[What the theorem does and does not show]
\label{rem:entropy_limits}
\begin{enumerate}[label=(\roman*)]
    \item The step realization is not invariant under simultaneous permutations of rows and columns, so it is to be expected that it separates some isospectral pairs. The content of Theorem~\ref{thm:isospectral_separation} is the explicit computation.
    \item Persistent entropy depends only on the multiset of lifetimes, so it cannot separate two matrices whose diagrams have the same lifetimes. For example, let $J$ be the reversal permutation $i \mapsto 6 - i$ and $C = JB'J^T$. Then $C \ne B'$, $C$ is isospectral to $B'$, and $\Phi_\square(C)$ is $\Phi_\square(B')$ rotated by $\pi$ about the centre of the square. Hence $\mathrm{Dgm}_1(\Phi_\square(C)) = \mathrm{Dgm}_1(\Phi_\square(B'))$, which has two bars of different lifetimes, and $E^{(1)}_{\mathrm{pers}}(C) = E^{(1)}_{\mathrm{pers}}(B')$. So two bars of different lifetimes do not force the entropies of an isospectral pair to differ.
    \item In the $3 \times 3$ example, perturbing the eight outer entries to distinct values does not create a second $H_1$ bar: a $3 \times 3$ grid has only one interior cell, so $\beta_1(X_t) \le 1$ for every $t$.
\end{enumerate}
The diagrams and entropies above were also checked numerically by two independent methods: union--find on the cells, using Alexander duality, and pixel rasterization with connected-component counting.
\end{remark}

% =========================================================================
\section{Level-Set Curvature Energies and Isoperimetry}
\label{sec:willmore_gmt}
% =========================================================================

\subsection{Curvature of Level Sets}
Let $\Omega \subset \R^d$ be a bounded open set, $d \ge 2$, and $\Phi(A) \in C^2(\overline{\Omega})$. On the open set $\{\nabla \Phi(A) \ne 0\}$ the level sets
\begin{equation}
    \Sigma_t \coloneqq \left\{ x \in \Omega \;\middle|\; \Phi(A)(x) = t \right\}
\end{equation}
are $C^2$ hypersurfaces with unit normal
\begin{equation}
    \mathbf{n}(x) \coloneqq \frac{\nabla \Phi(A)(x)}{\|\nabla \Phi(A)(x)\|}.
\end{equation}
If $\Phi(A) \in C^d$, Sard's theorem \cite{sard1942measure} shows that almost every $t$ is a regular value.

\begin{definition}[Level-set curvature]
On $\{\nabla \Phi(A) \ne 0\}$ let
\begin{equation}
    H_t(x) \coloneqq -\operatorname{div}\left( \frac{\nabla \Phi(A)(x)}{\|\nabla \Phi(A)(x)\|} \right) = -\frac{\Delta \Phi(A)}{\|\nabla \Phi(A)\|} + \frac{\Hess \Phi(A)(\nabla \Phi(A), \nabla \Phi(A))}{\|\nabla \Phi(A)\|^3}.
\end{equation}
Up to sign, $H_t(x)$ is the sum $\kappa_1 + \dots + \kappa_{d-1}$ of the principal curvatures of $\Sigma_t$ at $x$, that is, $(d-1)$ times the mean curvature.
\end{definition}

\subsection{The Integrated Willmore Energy}
The Willmore energy of a closed surface $\Sigma \subset \R^3$ is $\int_\Sigma H_{\mathrm{m}}^2 \dif A$, where $H_{\mathrm{m}} = (\kappa_1 + \kappa_2)/2$ is the mean curvature; Marques and Neves proved Willmore's conjecture that it is at least $2\pi^2$ for every torus immersed in $\R^3$ \cite{marques2014minmax}. We integrate the analogous quantity over all level sets.

\begin{definition}[Integrated Willmore energy]
For $d \ge 2$ and $\Phi(A) \in C^2(\overline{\Omega})$ whose critical set has measure zero,
\begin{equation}
    \mathcal{W}(\Phi(A)) \coloneqq \int_{-\infty}^\infty \left( \int_{\Sigma_t \cap \{\nabla\Phi(A) \ne 0\}} H_t(x)^2 \dif \mathcal{H}^{d-1}(x) \right) \dif t \in [0, \infty].
\end{equation}
By the coarea formula for Lipschitz functions \cite{evans2015measure}, this equals
\begin{equation}
    \mathcal{W}(\Phi(A)) = \int_{\{\nabla\Phi(A) \ne 0\}} \left| \operatorname{div}\left(\frac{\nabla \Phi(A)}{\|\nabla \Phi(A)\|}\right) \right|^2 \|\nabla \Phi(A)\| \dif x.
    \label{eq:willmore_volume}
\end{equation}
\end{definition}
Functionals of the form $\int |\nabla u|\,(\alpha + \beta|\operatorname{div}(\nabla u/|\nabla u|)|^p)\dif x$ were proposed for image inpainting; Chan, Kang and Shen study the case $p = 2$ in the plane (Euler's elastica) \cite{chan2003euler}, and Ambrosio and Masnou study the relaxation of the general case \cite{ambrosio2003direct}. The right-hand side of \eqref{eq:willmore_volume} is the term with $\alpha = 0$, $\beta = 1$, $p = 2$.

\begin{remark}[Integrability]
The formula for $H_t$ gives, with the Frobenius norm of the Hessian,
\begin{equation}
    |H_t| \le (\sqrt{d} + 1)\frac{\|\Hess \Phi(A)\|}{\|\nabla \Phi(A)\|}, \qquad H_t^2\|\nabla\Phi(A)\| \le (\sqrt{d}+1)^2 \frac{\|\Hess\Phi(A)\|^2}{\|\nabla \Phi(A)\|}.
\end{equation} Near a non-degenerate critical point $x_0$ we have $\|\nabla \Phi(A)(x)\| \ge c\|x - x_0\|$, so the integrand is $O(\|x - x_0\|^{-1})$, which is integrable in dimension $d \ge 2$. Hence $\mathcal{W}(\Phi(A)) < \infty$ when $\Phi(A)$ is a Morse function on a neighbourhood of $\overline{\Omega}$. In dimension $d = 1$ level sets are points and $H_t \equiv 0$, so the definition is used only for $d \ge 2$.
\end{remark}

\begin{proposition}[Conformal invariance; classical]
\label{thm:willmore_invariance}
Let $d = 3$ and let $\Sigma = \Sigma_t \subset \operatorname{int}\Omega$ be a compact regular level surface of $\Phi(A)$, with principal curvatures $\kappa_1, \kappa_2$, $H = \kappa_1 + \kappa_2$ (Definition above, up to sign) and Gaussian curvature $K = \kappa_1\kappa_2$. Let $\psi$ be a M\"obius transformation of $\R^3 \cup \{\infty\}$ whose pole, if any, does not lie on $\Sigma$. Then
\begin{equation}
    \int_{\psi(\Sigma)} \Big(\frac{H^2}{4} - K\Big) \dif \mathcal{H}^2 = \int_{\Sigma} \Big(\frac{H^2}{4} - K\Big) \dif \mathcal{H}^2
    \quad\text{and}\quad
    \int_{\psi(\Sigma)} (H^2 - K) \dif \mathcal{H}^2 = \int_{\Sigma} (H^2 - K) \dif \mathcal{H}^2.
\end{equation}
\end{proposition}

\begin{proof}
The first identity is White's theorem \cite{white1973global}: with the mean curvature $H_{\mathrm{m}} = H/2$, the integrand $H_{\mathrm{m}}^2 - K = \frac14(\kappa_1 - \kappa_2)^2$ times the area element is invariant under M\"obius transformations. For the second, $H^2 - K = 4(H^2/4 - K) + 3K$, and $\int_\Sigma K \dif\mathcal{H}^2 = 2\pi\chi(\Sigma)$ by Gauss--Bonnet; $\psi(\Sigma)$ is a closed surface diffeomorphic to $\Sigma$, so this term is also unchanged. Note that $\psi(\Sigma)$ is in general not a level set of $\Phi(A)$.
\end{proof}

\begin{proposition}[Willmore energy of a trigonometric checkerboard]
\label{prop:willmore_scaling}
Let $d = 2$, $\Omega = (0,1)^2$, $u(y) = \sin y_1 \sin y_2$, and for $k \ge 1$, $\epsilon > 0$ let $\Phi_k(x) = \epsilon\, u(kx)$. Put
\begin{equation}
    c_* \coloneqq \frac{1}{(2\pi)^2} \int_{[0,2\pi]^2} H_u(y)^2 \|\nabla u(y)\| \dif y,
\end{equation}
where $H_u$ is the level-set curvature of $u$. Then $0 < c_* < \infty$, and
\begin{equation}
    \int_\Omega \|\nabla \Phi_k\|^2 \dif x = \frac{\epsilon^2 k^2}{2}\big(1 + O(k^{-1})\big), \qquad \mathcal{W}(\Phi_k) = c_*\, \epsilon k^3 \big(1 + O(k^{-1})\big),
\end{equation}
with $O(k^{-1})$ uniform in $\epsilon$. In particular, for $\epsilon = 1/k$ the Dirichlet energy tends to $\frac12$ while $\mathcal{W}(\Phi_k) \sim c_* k^2$; for fixed $\epsilon$, $\mathcal{W}(\Phi_k) \sim c_* \epsilon k^3$.
\end{proposition}

\begin{proof}
The function $u$ is Morse. Its critical points are the zeros of both $\cos y_1 \sin y_2$ and $\sin y_1 \cos y_2$, and there the Hessian is $\pm\begin{psmallmatrix} 0 & 1 \\ 1 & 0 \end{psmallmatrix}$ or $\pm\Id$. By the integrability remark above, $g \coloneqq H_u^2\|\nabla u\|$ is locally integrable, so $c_* < \infty$. Near a strict maximum the level curves of $u$ are closed convex curves with non-zero curvature, so $c_* > 0$. The level-set curvature does not change when a function is multiplied by $\epsilon > 0$, and it scales like the inverse of a length under $x \mapsto kx$, so $H_{\Phi_k}(x) = k H_u(kx)$ and $\|\nabla\Phi_k(x)\| = \epsilon k \|\nabla u(kx)\|$. Substituting $y = kx$ in \eqref{eq:willmore_volume},
\begin{equation}
    \mathcal{W}(\Phi_k) = \epsilon k^3 \cdot \frac{1}{k^2} \int_{(0,k)^2} g(y) \dif y.
\end{equation}
The function $g \ge 0$ is $2\pi$-periodic in each variable. Let $m = \lfloor k / 2\pi \rfloor$. The square $(0,k)^2$ contains $m^2$ period cells and is covered by $(m+1)^2$ of them, so $(2\pi m/k)^2 c_* \le k^{-2}\int_{(0,k)^2} g \le (2\pi(m+1)/k)^2 c_*$, and both bounds equal $c_*(1 + O(k^{-1}))$. The same argument applied to $\|\nabla u\|^2 = \cos^2 y_1 \sin^2 y_2 + \sin^2 y_1 \cos^2 y_2$, whose mean over a period is $\frac12$, gives the Dirichlet energy.
\end{proof}

\begin{remark}
Numerically, $c_* \approx 1.058$; midpoint quadrature over one period with $1000^2$, $2000^2$ and $4000^2$ points gives $1.0578$ in each case, and direct quadrature of $\mathcal{W}(\Phi_k)/(\epsilon k^3)$ on $(0,1)^2$ gives $1.084$, $1.062$, $1.058$ for $k = 16, 32, 64$. The profile $\Phi_k$ is a model for a $k \times k$ checkerboard matrix $A_k$ with entries $\pm\epsilon$, for which $\|A_k\|_F = \epsilon k$; it is not produced by any realization operator defined in this paper. The step realization of $A_k$ is not $C^2$, and $\mathcal{W}$ is not defined for it.
\end{remark}

\subsection{The Isoperimetric Profile}
\begin{definition}[Isoperimetric profile and Cheeger constant]
Let $\Omega = (0,1)^d$ and let $\Phi(A) \in C^1(\overline{\Omega})$ be a density realization with $\Phi(A) \ge c > 0$, inducing $\mu_A$. For a set $E \subset \Omega$ of finite perimeter, let $\Per_{\mu_A}(E) \coloneqq \int_{\partial^* E \cap \Omega} \Phi(A) \dif \mathcal{H}^{d-1}$ be its weighted perimeter relative to $\Omega$. The isoperimetric profile $\mathcal{I}_A \colon (0, 1) \to [0,\infty)$ and the Cheeger constant are
\begin{equation}
    \mathcal{I}_A(v) \coloneqq \inf \left\{ \Per_{\mu_A}(E) \;\middle|\; \mu_A(E) = v \right\}, \qquad h(A) \coloneqq \inf_{0 < \mu_A(E) \le 1/2} \frac{\Per_{\mu_A}(E)}{\mu_A(E)}.
\end{equation}
\end{definition}

\begin{proposition}[Cheeger's inequality; classical]
\label{prop:cheeger}
Let $\lambda_1(A) \coloneqq \inf\{ \int_\Omega \|\nabla f\|^2 \dif\mu_A / \int_\Omega (f - \bar f)^2 \dif\mu_A \}$ over non-constant $f \in H^1(\Omega)$, with $\bar f = \int f \dif\mu_A$; this is the first non-zero eigenvalue of $-\Delta_A$ with Neumann conditions. Then $\lambda_1(A) \ge h(A)^2/4$.
\end{proposition}

\begin{proof}
This is Cheeger's argument \cite{cheeger1970lower}. By density it suffices to take $f \in C^1(\overline{\Omega})$ non-constant. Let $m$ be a median of $f$ for $\mu_A$ and $g_\pm \coloneqq (f - m)_\pm$. For $t > 0$ the sets $\{g_+^2 > t\}$ have $\mu_A$-measure at most $\frac12$. The coarea formula for the Lipschitz function $g_+^2$ \cite{evans2015measure} gives
\begin{equation}
    \int_\Omega \|\nabla g_+^2\| \dif\mu_A = \int_0^\infty \Per_{\mu_A}(\{g_+^2 > t\}) \dif t \ge h(A) \int_0^\infty \mu_A(\{g_+^2 > t\}) \dif t = h(A) \int_\Omega g_+^2 \dif\mu_A.
\end{equation}
By Cauchy--Schwarz, $\int \|\nabla g_+^2\| \dif\mu_A = 2\int g_+\|\nabla g_+\| \dif\mu_A \le 2 (\int g_+^2 \dif\mu_A)^{1/2} (\int \|\nabla g_+\|^2 \dif\mu_A)^{1/2}$, so $\int \|\nabla g_+\|^2 \dif\mu_A \ge \frac{h(A)^2}{4}\int g_+^2 \dif\mu_A$, and likewise for $g_-$. Adding, and using that $\nabla g_+$ and $\nabla g_-$ have disjoint supports up to null sets, $\int \|\nabla f\|^2 \dif\mu_A \ge \frac{h(A)^2}{4} \int (f-m)^2 \dif\mu_A \ge \frac{h(A)^2}{4} \int (f - \bar f)^2 \dif\mu_A$.
\end{proof}

\begin{remark}
On compact Riemannian manifolds without boundary whose Ricci curvature is bounded below, an inequality in the opposite direction bounds $\lambda_1$ above in terms of the Cheeger constant; it is due to Buser \cite{buser1982note}, and Ledoux gave an analytic proof \cite{ledoux1994simple}. We do not use it.
\end{remark}

% =========================================================================
\section{Wavefront Sets and Besov Regularity}
\label{sec:microlocal}
% =========================================================================

\subsection{Gabor Transform}
To record the frequency and direction $\xi \in \R^d \setminus \{0\}$ of features of $\Phi(A)$ as well as their position, we work on $T^*\Omega \cong \Omega \times \R^d$. Let $g(y) = (\pi s^2)^{-d/4} e^{-\|y\|^2/(2s^2)}$, $s > 0$, be a Gaussian window with $\|g\|_{L^2} = 1$, and extend $\Phi(A)$ by zero outside $\Omega$.
\begin{definition}[Gabor transform]
For $x \in \Omega$ and $\xi \in \R^d$, the Gabor transform (short-time Fourier transform with Gaussian window \cite{grochenig2001foundations}) of $\Phi(A) \in L^2(\Omega)$ is:
\begin{equation}
    \mathcal{V}_{\Phi(A)}(x, \xi) \coloneqq \int_{\R^d} \Phi(A)(y) g(y - x) e^{-i \langle \xi, y \rangle} \dif y.
\end{equation}
The associated phase-space energy density (spectrogram) is:
\begin{equation}
    \mathcal{E}_A(x, \xi) \coloneqq \left| \mathcal{V}_{\Phi(A)}(x, \xi) \right|^2 \ge 0.
\end{equation}
\end{definition}

\subsection{The Wavefront Set}
We use H\"ormander's definition \cite[Sect.~8.1]{hormander2003analysis}.

\begin{definition}[Wavefront set]
Let $f \in L^1_{\mathrm{loc}}(\Omega)$. A point $(x_0, \xi_0) \in \Omega \times (\R^d \setminus \{0\})$ is \emph{not} in $\WF(f)$ if there are $\chi \in C_c^\infty(\Omega)$ with $\chi(x_0) \ne 0$ and an open cone $\Gamma \ni \xi_0$ such that, for every $N \in \mathbb{N}$,
\begin{equation}
    \sup_{\xi \in \Gamma} (1 + \|\xi\|)^N \big| \widehat{\chi f}(\xi) \big| < \infty.
\end{equation}
\end{definition}

\begin{remark}[A fixed window does not localize]
\label{rem:fixed_window}
The same decay condition imposed on the Gabor transform with a fixed Gaussian window does not localize singularities. For $d = 1$ and $f = \mathbf{1}_{[0,\infty)}$, integration by parts gives, for every $x_0 \in \R$,
\begin{equation}
    \mathcal{V}_f(x_0, \xi) = \int_0^\infty g(y - x_0) e^{-i\xi y} \dif y = \frac{g(-x_0)}{i\xi} + \frac{1}{i\xi}\int_0^\infty g'(y - x_0) e^{-i\xi y}\dif y = \frac{g(-x_0)}{i\xi} + O(\xi^{-2}),
\end{equation}
and $g(-x_0) \ne 0$. So $|\mathcal{V}_f(x_0,\xi)|$ decays only like $|\xi|^{-1}$ even at points $x_0$ far from the jump, because the Gaussian tail reaches the jump. Numerically, at $x_0 = 5s$ one finds $|\xi\, \mathcal{V}_f(x_0,\xi)| \to g(-x_0)$. For this reason the wavefront set is defined with compactly supported cut-offs.
\end{remark}

\begin{proposition}[Jump across a smooth interface; classical]
\label{thm:wavefront_interface}
Let $U \subset \Omega$ be open, $D \subset \R^d$ open with $C^\infty$ boundary, $\Gamma = \partial D \cap U$, and let $f = b + (a - b)\mathbf{1}_D$ on $U$ with constants $a \ne b$. Then
\begin{equation}
    \WF(f) \cap (U \times \R^d) = N^*\Gamma \setminus 0 \coloneqq \left\{ (x, \lambda \mathbf{n}_\Gamma(x)) \;\middle|\; x \in \Gamma, \; \lambda \ne 0 \right\}.
\end{equation}
\end{proposition}

\begin{proof}
Off $\Gamma$, $f$ is locally constant, hence smooth, so no point over $U \setminus \Gamma$ lies in $\WF(f)$. Let $x_0 \in \Gamma$. By the transformation law of wavefront sets under diffeomorphisms \cite[Sect.~8.2]{hormander2003analysis}, which maps conormal bundles to conormal bundles, we may assume near $x_0$ that $\Gamma = \{y_1 = 0\}$ and $f = b + (a-b)H(y_1)$, $H = \mathbf{1}_{(0,\infty)}$. Let $\chi \in C_c^\infty$ near $x_0$ and write $\xi = (\xi_1, \xi')$. Since $\partial_{y_j}$, $j \ge 2$, does not see the jump, $(\xi')^{\alpha}\, \widehat{\chi H}(\xi) = \widehat{(-i\partial')^{\alpha}\chi \cdot H}(\xi)$ up to sign, which is bounded by $\|\partial'^{\alpha}\chi\|_{L^1}$. In a closed cone not containing $\pm e_1$ we have $\|\xi'\| \ge c\|\xi\|$, so $\widehat{\chi f}$ decays rapidly there. Hence $\WF(f)$ over $x_0$ is contained in $\{\lambda e_1 : \lambda \ne 0\}$. It is non-empty, because $x_0$ lies in the singular support of $f$, which is the projection of $\WF(f)$ \cite[Sect.~8.1]{hormander2003analysis}. Since $f$ is real and we may take $\chi$ real, $\widehat{\chi f}(-\xi) = \overline{\widehat{\chi f}(\xi)}$, so $\WF(f)$ is invariant under $\xi \mapsto -\xi$. Therefore $\WF(f)$ over $x_0$ equals $\{\lambda e_1 : \lambda \ne 0\}$.
\end{proof}

\begin{remark}[Step realizations]
For the step realization of Definition~\ref{def:step}, Proposition~\ref{thm:wavefront_interface} applies near every point in the relative interior of a cell edge separating two different values, and $\WF$ there consists of the normals to the edge. At a grid vertex the set can be larger. For instance, for $f = H(y_1)H(y_2)$ and $\chi = \beta(y_1)\beta(y_2)$ with $0 \le \beta \in C_c^\infty$, $\beta(0) > 0$, one has $\widehat{\chi f}(\xi) = \widehat{\beta H}(\xi_1)\widehat{\beta H}(\xi_2)$ with $\widehat{\beta H}(t) = \beta(0)/(it) + O(t^{-2})$. This decays only polynomially in every direction, for cut-offs of arbitrarily small support, so every direction belongs to $\WF(f)$ over the vertex.
\end{remark}

\subsection{Besov Regularity Index}
Let $\Omega = (0,1)^d$ and $\Omega_h \coloneqq \{x \in \Omega : x + h \in \Omega\}$. For $1 \le p < \infty$ and $s \in (0,1)$ we use the first-difference seminorm of the Besov space $B^s_{p,\infty}$ (see \cite{triebel1983theory} for the Besov scale)
\begin{equation}
    [\Phi(A)]_{B^s_{p,\infty}} \coloneqq \sup_{h \in \R^d, 0 < \|h\| \le 1} \frac{\|\Phi(A)(\cdot + h) - \Phi(A)(\cdot)\|_{L^p(\Omega_h)}}{\|h\|^s}.
\end{equation}

\begin{definition}[Besov regularity index]
For $1 \le p < \infty$,
\begin{equation}
    s^*_p(A) \coloneqq \sup \left\{ s \in (0, 1) \;\middle|\; [\Phi(A)]_{B^s_{p,\infty}} < \infty \right\} \in [0,1].
\end{equation}
\end{definition}

\begin{proposition}[Besov index of smooth and step realizations]
\label{prop:besov}
Let $1 \le p < \infty$.
\begin{enumerate}[label=(\alph*)]
    \item If $\Phi(A) \in C^1(\overline{\Omega})$, then $s^*_p(A) = 1$.
    \item Let $d = 2$ and let $\Phi_\square(A)$ be the step realization of Definition~\ref{def:step}, with $\Phi_\square(A)$ not constant. Then $s^*_p(A) = 1/p$ for $p > 1$, and $s^*_1(A) = 1$.
\end{enumerate}
\end{proposition}

\begin{proof}
(a) $\Omega$ is convex, so $|f(x+h) - f(x)| \le \|\nabla f\|_\infty \|h\|$ and $[f]_{B^s_{p,\infty}} \le \|\nabla f\|_\infty < \infty$ for every $s < 1$.

(b) Write $f = \Phi_\square(A)$. If $x, x + h \in \Omega$ and $f(x+h) \ne f(x)$ (away from the null set of grid lines), the segment from $x$ to $x+h$ crosses one of the $2(n-1)$ interior grid lines $\{x_i = j/n\}$. For a fixed line $\{x_i = c\}$ this happens only for $x$ in a slab of width $|h_i| \le \|h\|$. Hence $\|f(\cdot + h) - f\|_{L^p(\Omega_h)}^p \le (2\|f\|_\infty)^p\, 2(n-1)\|h\|$, so $[f]_{B^s_{p,\infty}} < \infty$ for $s \le 1/p$ and $s^*_p \ge \min(1, 1/p)$; this gives $s^*_1 = 1$. Conversely, since $f$ is not constant, two cells sharing an edge carry different values $\alpha \ne \beta$; say the edge lies on $\{x_i = c\}$, with the cell of value $\alpha$ below it. For $h = t e_i$, $0 < t < 1/n$, every $x$ in the open lower cell with $x_i > c - t$ satisfies $x + h$ in the open upper cell. This set has measure $t/n$, so $\|f(\cdot+h) - f\|_{L^p(\Omega_h)} \ge |\alpha - \beta|\, (t/n)^{1/p}$, and $[f]_{B^s_{p,\infty}} = \infty$ for every $s > 1/p$. Thus $s^*_p = 1/p$ when $p > 1$.
\end{proof}

\begin{remark}
For $p = 1$ the index does not distinguish step realizations from smooth ones, so $p > 1$ must be used for that purpose. For an indicator $f = \mathbf{1}_E$ of a measurable set $E \subset \Omega$, a short argument gives a lower bound: if exactly one of $x, x+h$ lies in $E$, the segment between them meets $\partial E$, so $\|f(\cdot + h) - f\|_{L^1(\Omega_h)}$ is at most the measure of the $\|h\|$-neighbourhood of $\partial E$. By the definition of the upper Minkowski dimension $\overline{\dim}_M$, this yields $s^*_1 \ge \min\{1, d - \overline{\dim}_M \partial E\}$. We do not claim equality. For the fractional seminorm $[\mathbf{1}_E]_{W^{s,1}}$ (the fractional perimeter) Visintin proved the analogous bound, finiteness for $s < d - \overline{\dim}_M \partial E$ \cite{visintin1991generalized}, and Lombardini showed that the resulting index equals the Minkowski codimension for the von Koch snowflake \cite{lombardini2019fractional}.
\end{remark}

% =========================================================================
\section{Dixmier Traces and Cyclic Cocycles}
\label{sec:noncommutative}
% =========================================================================

\subsection{The Spectral Triple on the Torus}
We use Connes' non-commutative geometry \cite{connes1994noncommutative}. In this section $\Omega = \mathbb{T}^d = (\R/\Z)^d$ with the flat metric and the trivial spin structure, and the realizations $\Phi(A)$ are real and smooth on $\mathbb{T}^d$ (for the kernel realization, take periodic $\psi_i \in C^\infty(\mathbb{T}^d)$). Let $N_d = 2^{\lfloor d/2 \rfloor}$ and let $\gamma^1, \dots, \gamma^d$ be Hermitian $N_d \times N_d$ matrices with $\gamma^\mu\gamma^\nu + \gamma^\nu\gamma^\mu = 2\delta^{\mu\nu}\Id$. We consider
\begin{itemize}
    \item $\mathcal{H} \coloneqq L^2(\mathbb{T}^d, \C^{N_d})$;
    \item $\mathcal{D} \coloneqq i \gamma^\mu \partial_\mu$, self-adjoint with discrete spectrum. Its kernel consists of the constant spinors and has dimension $N_d$. We set $|\mathcal{D}|^{-d} \coloneqq 0$ on $\ker \mathcal{D}$; on the Fourier mode $e^{2\pi i \langle k, x\rangle}$, $k \in \Z^d \setminus \{0\}$, it acts as $|2\pi k|^{-d}$;
    \item the multiplication operators $M_{\Phi(A)} \colon \psi \mapsto \Phi(A)\,\psi$.
\end{itemize}

\subsection{The Dixmier Trace of a Realization}
\begin{definition}[Dixmier trace]
For a positive compact operator $T$ in the ideal $\mathcal{L}^{1,\infty}(\mathcal{H})$, with singular values $\mu_1(T) \ge \mu_2(T) \ge \dots$, put
\begin{equation}
    \Tr_\omega(T) \coloneqq \omega\text{-}\lim_{N \to \infty} \frac{1}{\log N} \sum_{n=1}^N \mu_n(T),
\end{equation}
where $\omega$ is a Dixmier limit, and extend $\Tr_\omega$ by linearity to all of $\mathcal{L}^{1,\infty}(\mathcal{H})$ \cite[Ch.~IV, Sect.~2]{connes1994noncommutative}; see \cite{lord2013singular} for the general theory of singular traces. For the operators below the value does not depend on $\omega$.
\end{definition}

\begin{proposition}[Dixmier trace of a realization]
\label{thm:dixmier_integration}
For real $\Phi(A) \in C^\infty(\mathbb{T}^d)$,
\begin{equation}
    \Tr_\omega \left( M_{\Phi(A)} |\mathcal{D}|^{-d} \right) = c_d \int_{\mathbb{T}^d} \Phi(A)(x) \dif x, \qquad c_d \coloneqq \frac{2^{\lfloor d/2 \rfloor} |\mathbb{S}^{d-1}|}{d (2\pi)^d},
\end{equation}
where $|\mathbb{S}^{d-1}| = 2\pi^{d/2}/\Gamma(d/2)$.
\end{proposition}

\begin{proof}
$|\mathcal{D}|^{-d}$ is the Fourier multiplier with symbol $\|\xi\|^{-d}\Id$ for $\xi \ne 0$; modifying it near $\xi = 0$ changes the operator by a finite-rank one. So $T = M_{\Phi(A)}|\mathcal{D}|^{-d}$ is a classical pseudo-differential operator of order $-d$ on the closed manifold $\mathbb{T}^d$, with principal symbol $\sigma_{-d}(T)(x,\xi) = \Phi(A)(x)\|\xi\|^{-d}\Id$. Connes' trace theorem \cite{connes1988action}, \cite[Ch.~IV, Sect.~2, Prop.~5]{connes1994noncommutative} states that such an operator lies in $\mathcal{L}^{1,\infty}$, that its Dixmier trace does not depend on $\omega$, and that it equals the noncommutative residue of Wodzicki and Guillemin \cite{wodzicki1987noncommutative,guillemin1985new}:
\begin{equation}
    \Tr_\omega(T) = \frac{1}{d (2\pi)^d} \int_{S^*\mathbb{T}^d} \operatorname{tr} \sigma_{-d}(T)(x, \xi) \dif S(x, \xi).
\end{equation}
Since $\operatorname{tr}\Id = 2^{\lfloor d/2 \rfloor}$, integration over the unit sphere gives $|\mathbb{S}^{d-1}|$, which proves the formula.
\end{proof}

\begin{remark}[Linearity matters]
For $T \ge 0$ the singular-value formula applies directly. For a sign-changing $\Phi(A)$, the operator $M_{\Phi(A)}|\mathcal{D}|^{-d}$ is not positive, and the plain sum of its singular values computes the Dixmier trace of $|T|$, not of $T$. Numerically, for $d = 1$ and $\Phi = \cos 2\pi x$, the singular-value sums grow like $0.202\log N \approx \frac{1}{\pi}\int|\Phi| \log N$, while the linear Dixmier trace (diagonal sums in the Fourier basis) is $\frac{1}{\pi}\int\Phi = 0$, in agreement with Proposition~\ref{thm:dixmier_integration}.
\end{remark}

\subsection{Cyclic Cocycles in Dimension Two}
The commutator of $\mathcal{D}$ with a realization is the multiplication operator
\begin{equation}
    [\mathcal{D}, M_{\Phi(A)}] = i \gamma^\mu \, \partial_\mu \Phi(A).
\end{equation}
For a tuple of $d+1$ matrices $(A_0, A_1, \dots, A_d)$ we define
\begin{equation}
    \tau_d(A_0, A_1, \dots, A_d) \coloneqq \Tr_\omega \left( \Phi(A_0) [\mathcal{D}, \Phi(A_1)] [\mathcal{D}, \Phi(A_2)] \cdots [\mathcal{D}, \Phi(A_d)] |\mathcal{D}|^{-d} \right).
\end{equation}
For $d = 2$ take $\gamma^1 = \sigma_1$, $\gamma^2 = \sigma_2$ (Pauli matrices) and the grading $\gamma \coloneqq -i\gamma^1\gamma^2 = \sigma_3$, which anticommutes with $\mathcal{D}$, and define the graded version
\begin{equation}
    \tau_2^\gamma(A_0, A_1, A_2) \coloneqq \Tr_\omega \left( \gamma\, \Phi(A_0) [\mathcal{D}, \Phi(A_1)] [\mathcal{D}, \Phi(A_2)] |\mathcal{D}|^{-2} \right).
\end{equation}

Up to a normalizing constant, $\tau_d$ and $\tau_2^\gamma$ are the Dixmier-trace cochains that Connes uses to compute the Hochschild class of the Chern character of a spectral triple, with $\gamma\Phi(A_0)$ in place of $\Phi(A_0)$ in the even case \cite[Ch.~IV, Sect.~2, Thm.~8]{connes1994noncommutative}. Cyclic cohomology was introduced in \cite{connes1985noncommutative}, and the local index formula of Connes and Moscovici expresses the Chern character of a spectral triple through Dixmier traces and residues \cite{connes1995local}.

\begin{proposition}[Cocycles on $\mathbb{T}^2$]
\label{prop:cocycle}
Let $d = 2$, write $\Phi_j = \Phi(A_j)$, and orient $\mathbb{T}^2$ by $\dif x_1 \wedge \dif x_2$. Then:
\begin{enumerate}[label=(\alph*)]
    \item $\displaystyle \tau_2(A_0, A_1, A_2) = -\frac{1}{2\pi} \int_{\mathbb{T}^2} \Phi_0 \, \nabla\Phi_1 \cdot \nabla\Phi_2 \dif x$;
    \item $\displaystyle \tau_2^\gamma(A_0, A_1, A_2) = -\frac{i}{2\pi} \int_{\mathbb{T}^2} \Phi_0 \, \dif\Phi_1 \wedge \dif\Phi_2$;
    \item if $\Phi_0^2 + \Phi_1^2 + \Phi_2^2 = 1$, so that $F = (\Phi_0, \Phi_1, \Phi_2) \colon \mathbb{T}^2 \to \mathbb{S}^2$, then
    \begin{equation}
        \frac{i}{4} \sum_{\sigma \in S_3} \operatorname{sgn}(\sigma)\, \tau_2^\gamma(A_{\sigma(0)}, A_{\sigma(1)}, A_{\sigma(2)}) = \deg F \in \Z.
    \end{equation}
\end{enumerate}
\end{proposition}

\begin{proof}
The operator in (a) is $M_G|\mathcal{D}|^{-2}$ with the smooth matrix-valued function
\begin{equation}
    G = \Phi_0 (i\gamma^\mu\partial_\mu\Phi_1)(i\gamma^\nu\partial_\nu\Phi_2) = -\Phi_0\, \gamma^\mu\gamma^\nu\, \partial_\mu\Phi_1 \partial_\nu\Phi_2.
\end{equation}
As in the proof of Proposition~\ref{thm:dixmier_integration}, Connes' trace theorem gives
\begin{equation}
    \Tr_\omega(M_G|\mathcal{D}|^{-2}) = \frac{1}{2(2\pi)^2} \cdot 2\pi \int \operatorname{tr} G \dif x = \frac{1}{4\pi}\int \operatorname{tr}G \dif x.
\end{equation} Since $\operatorname{tr}(\gamma^\mu\gamma^\nu) = 2\delta^{\mu\nu}$, $\operatorname{tr} G = -2\Phi_0 \nabla\Phi_1\cdot\nabla\Phi_2$, which gives (a). For (b) the function is $\gamma G$. Since $\gamma = \sigma_3$ and $\sigma_\mu\sigma_\nu = \delta_{\mu\nu}\Id + i\epsilon_{\mu\nu 3}\sigma_3$, we get $\operatorname{tr}(\gamma\gamma^\mu\gamma^\nu) = 2i\epsilon^{\mu\nu}$, hence $\operatorname{tr}(\gamma G) = -2i\,\Phi_0(\partial_1\Phi_1\partial_2\Phi_2 - \partial_2\Phi_1\partial_1\Phi_2)$, which gives (b). For (c), $\sum_\sigma \operatorname{sgn}\sigma\, F^{\sigma(0)} \dif F^{\sigma(1)} \wedge \dif F^{\sigma(2)} = 2\, F\cdot(\partial_1 F \times \partial_2 F) \dif x_1\wedge\dif x_2 = 2F^*\omega$, where $\omega$ is the area form of $\mathbb{S}^2$ with the outward orientation. Since $\int_{\mathbb{T}^2} F^*\omega = 4\pi\deg F$, (b) gives $\sum_\sigma \operatorname{sgn}\sigma\, \tau_2^\gamma(\dots) = -\frac{i}{2\pi}\cdot 2\cdot 4\pi\deg F = -4i\deg F$.
\end{proof}

\begin{remark}[The ungraded cocycle is not topological]
\label{rem:tau_not_topological}
By (a), $\tau_2$ is symmetric in its last two arguments and has no antisymmetric part, so it cannot detect the degree. For example, let $0 < s < 1$, $c(x) = \cos 2\pi x_1$, $\Phi_1 = \Phi_2 = s c/\sqrt{2}$ and $\Phi_0 = \sqrt{1 - s^2c^2}$. Then $F$ takes values in $\mathbb{S}^2$, has degree $0$ (since $\partial_2 F = 0$), and $\tau_2 = -\frac{1}{2\pi}\int\Phi_0\|\nabla\Phi_1\|^2 \dif x < 0$. Likewise, composing a degree-one map $\mathbb{T}^2 \to \mathbb{S}^2$ with rotations of $\mathbb{S}^2$ keeps the degree fixed but changes $\tau_2$; numerically, $\tau_2$ ranges over about $[-0.015, 0.037]$ for five random rotations, while $\frac{i}{4}\sum_\sigma\operatorname{sgn}\sigma\,\tau_2^\gamma = 1.000$ in all five cases. We treat only $d = 2$; other dimensions are not considered here.
\end{remark}

% =========================================================================
\section{Quantum Fisher Information of Fields of Density Matrices}
\label{sec:quantum_info}
% =========================================================================

\subsection{Fields of Density Matrices}
Let $\Omega \subset \R^d$ be a bounded open set and let $x \mapsto \rho(x)$ be a $C^1$ map from a neighbourhood of $\overline{\Omega}$ into the $\chi \times \chi$ density matrices ($\rho(x) \succeq 0$, $\Tr\rho(x) = 1$). In other words, $\rho$ is a section of the trivial bundle $\Omega \times \mathrm{End}(\C^\chi)$. Such fields arise, for instance, as auxiliary (bond) density matrices of continuous matrix product states \cite{verstraete2010continuous}, where to each $x$ are attached a bond space $\C^\chi$ and operators $Q(x), R(x) \in \mathrm{End}(\C^\chi)$. Nothing below uses this structure.

\subsection{The Integrated Quantum Fisher Information Volume}
For $\rho(x) \succ 0$, the quantum Fisher information metric is
\begin{equation}
    g^{\mathrm{QFI}}_{\mu\nu}(x) \coloneqq \frac{1}{2} \Tr_\chi \left( \rho(x) \{ L_\mu(x), L_\nu(x) \} \right),
\end{equation}
where $\{A, B\} = AB + BA$, and $L_\mu(x)$ is the Symmetric Logarithmic Derivative (SLD) defined implicitly by:
\begin{equation}
    \partial_\mu \rho(x) = \frac{1}{2} \left( \rho(x) L_\mu(x) + L_\mu(x) \rho(x) \right).
    \label{eq:sld_implicit}
\end{equation}
The symmetric logarithmic derivative goes back to Helstrom \cite{helstrom1967minimum}. Braunstein and Caves identified the corresponding metric with statistical distinguishability of states \cite{braunstein1994statistical}; it equals four times the Bures metric built from Uhlmann's fidelity \cite{uhlmann1976transition}, see \cite[Sect.~2.4.2]{liu2020quantum}, and it belongs to Petz's family of monotone metrics on density matrices \cite{petz1996monotone}. Pulling such metrics back to a parameter space is standard: for pure states this gives the metric of Provost and Vallee \cite{provost1980riemannian}, and the fidelity metric on the coupling space of a Hamiltonian is used to detect quantum phase transitions \cite{zanardi2007information}.

\begin{definition}[Integrated QFI volume]
The integrated quantum Fisher information volume of the field $\rho$ is
\begin{equation}
    \mathrm{Vol}_{\mathrm{QFI}}(\rho) \coloneqq \int_\Omega \sqrt{\det g^{\mathrm{QFI}}(x)} \dif^d x.
\end{equation}
\end{definition}

\begin{lemma}[QFI metric in an eigenbasis; classical, see {\cite[Thm.~2.1]{liu2020quantum}}]
\label{lem:qfi_formula}
Let $\rho(x) = \sum_i \lambda_i |i\rangle\langle i|$ with all $\lambda_i > 0$. Then $\langle i|L_\mu|j\rangle = 2\langle i|\partial_\mu\rho|j\rangle / (\lambda_i + \lambda_j)$ and
\begin{equation}
    g^{\mathrm{QFI}}_{\mu\nu}(x) = 2 \sum_{i,j=1}^\chi \frac{\operatorname{Re} \left( \langle i | \partial_\mu \rho | j \rangle \langle j | \partial_\nu \rho | i \rangle \right)}{\lambda_i + \lambda_j}.
    \label{eq:qfi_eigen}
\end{equation}
\end{lemma}

\begin{proof}
Taking the $(i,j)$ entry of \eqref{eq:sld_implicit} gives $\langle i|\partial_\mu\rho|j\rangle = \frac12(\lambda_i + \lambda_j)\langle i|L_\mu|j\rangle$. Then
\begin{equation}
    \frac12\Tr(\rho\{L_\mu,L_\nu\}) = \frac12\sum_{i,j}(\lambda_i + \lambda_j)\operatorname{Re}\big(\langle i|L_\mu|j\rangle\langle j|L_\nu|i\rangle\big),
\end{equation}
and substituting the entries of $L_\mu$, $L_\nu$ gives \eqref{eq:qfi_eigen}.
\end{proof}

\begin{theorem}[Vanishing of the QFI volume]
\label{thm:qfi_volume}
Assume $\rho(x) \succ 0$ for all $x \in \overline{\Omega}$. Then:
\begin{enumerate}[label=(\alph*)]
    \item $g^{\mathrm{QFI}}(x)$ is the Gram matrix of $\partial_1\rho(x), \dots, \partial_d\rho(x)$ for the real inner product
    \begin{equation}
        \langle X, Y\rangle_{\rho(x)} \coloneqq 2\sum_{i,j}\frac{\operatorname{Re}(X_{ij}\overline{Y_{ij}})}{\lambda_i + \lambda_j}
    \end{equation}
    on Hermitian matrices (entries in an eigenbasis of $\rho(x)$), which is positive definite. Hence $g^{\mathrm{QFI}}(x) \succeq 0$, $\mathrm{Vol}_{\mathrm{QFI}}(\rho) \ge 0$, and $g^{\mathrm{QFI}}_{\mu\mu}(x) \ge \|\partial_\mu\rho(x)\|_{\mathrm{HS}}^2$.
    \item $\det g^{\mathrm{QFI}}(x) > 0$ if and only if $\partial_1\rho(x), \dots, \partial_d\rho(x)$ are linearly independent over $\R$.
    \item $\mathrm{Vol}_{\mathrm{QFI}}(\rho) = 0$ if and only if the derivatives $\partial_1\rho(x), \dots, \partial_d\rho(x)$ are linearly dependent at every $x \in \Omega$, that is, $\rank \dif\rho(x) < d$ everywhere.
    \item $\mathrm{Vol}_{\mathrm{QFI}}$ is invariant under $C^1$ diffeomorphisms of $\Omega$: if $\phi \colon \Omega' \to \Omega$ is one, then $\mathrm{Vol}_{\mathrm{QFI}}(\rho \circ \phi) = \mathrm{Vol}_{\mathrm{QFI}}(\rho)$.
\end{enumerate}
\end{theorem}

\begin{proof}
(a) In the eigenbasis of $\rho(x)$ the matrices $\partial_\mu\rho$ are Hermitian, so $\langle j|\partial_\nu\rho|i\rangle = \overline{\langle i|\partial_\nu\rho|j\rangle}$, and \eqref{eq:qfi_eigen} reads $g_{\mu\nu} = \langle\partial_\mu\rho, \partial_\nu\rho\rangle_{\rho}$. The form $\langle\cdot,\cdot\rangle_\rho$ is a weighted Hilbert--Schmidt inner product with positive weights $2/(\lambda_i + \lambda_j)$; since $\lambda_i + \lambda_j \le 2$, the weights are at least $1$, which gives the bound on $g_{\mu\mu}$. A Gram matrix is positive semi-definite.
(b) A Gram matrix of $d$ vectors for a positive-definite inner product is non-singular if and only if the vectors are linearly independent.
(c) The SLD $L_\mu(x)$ is the unique solution of the linear equation \eqref{eq:sld_implicit}. The operator $X \mapsto \frac12(\rho X + X\rho)$ is invertible for $\rho \succ 0$ and depends continuously on $\rho$, so $g^{\mathrm{QFI}}$ and $\det g^{\mathrm{QFI}}$ are continuous on $\Omega$. Hence $\int_\Omega \sqrt{\det g^{\mathrm{QFI}}} = 0$ if and only if $\det g^{\mathrm{QFI}} \equiv 0$ on the open set $\Omega$, and (b) applies.
(d) By the chain rule, $g^{\mathrm{QFI}}_{\rho\circ\phi}(y) = J(y)^T g^{\mathrm{QFI}}_\rho(\phi(y)) J(y)$ with $J = \dif\phi$. So $\sqrt{\det g_{\rho\circ\phi}(y)} = \sqrt{\det g_\rho(\phi(y))}\,|\det J(y)|$, and the change-of-variables formula gives the claim.
\end{proof}

\begin{remark}[What the volume does and does not detect]
\label{rem:qfi_examples}
\begin{enumerate}[label=(\roman*)]
    \item Non-constant fields can have zero volume. For example, let $d = 2$ and let $\rho(x) = R(x_1 + x_2)$, where $R$ is a non-constant $C^1$ curve of positive definite density matrices. Then $\partial_1\rho = \partial_2\rho \ne 0$ and $\mathrm{Vol}_{\mathrm{QFI}} = 0$.
    \item A field that is constant up to a unitary gauge can have positive volume. If $\rho(x) = U(x)\rho_0U(x)^\dagger$ with $U(x) = \exp(i(x_1h_1 + x_2h_2))$, then $\partial_\mu\rho = i[h_\mu, \rho_0]$ at $x = 0$. For generic Hermitian $h_1, h_2$ and non-scalar $\rho_0$ these are linearly independent, so $\det g^{\mathrm{QFI}} > 0$ near $0$. Numerically, $\det g^{\mathrm{QFI}}(0) > 0$ in random $3 \times 3$ examples.
    \item Nothing in Theorem~\ref{thm:qfi_volume} relates the volume to the presence or absence of entanglement.
\end{enumerate}
Formula \eqref{eq:qfi_eigen} agrees with an independent computation of $L_\mu$ from \eqref{eq:sld_implicit} by a linear solve, in random examples; the same formula without the factor $2$ does not.
\end{remark}

\begin{remark}[Rank-deficient states]
When some eigenvalues of $\rho(x)$ vanish, the QFI metric can be discontinuous, as shown by {\v{S}}afr{\'a}nek \cite{safranek2017discontinuities}. The assumption $\rho \succ 0$ on $\overline{\Omega}$ excludes this.
\end{remark}

\subsection{An Entropy-Weighted Density}
Let $S_{\mathrm{vN}}(x) = -\Tr(\rho(x) \log \rho(x))$ be the von Neumann entropy of $\rho(x)$.

\begin{definition}[QFI-weighted entropy density and centroid]
Put
\begin{equation}
    s(x) \coloneqq S_{\mathrm{vN}}(x)\sqrt{\det g^{\mathrm{QFI}}(x)}, \qquad S_{\mathrm{tot}}(\rho) \coloneqq \int_\Omega s(x) \dif^d x.
\end{equation}
If $S_{\mathrm{tot}}(\rho) > 0$, the centroid and the dispersion of $s$ are
\begin{equation}
    \bar{x}_{s} \coloneqq \frac{1}{S_{\mathrm{tot}}(\rho)} \int_\Omega x \, s(x) \dif^d x, \quad \Sigma_{s}^2 \coloneqq \frac{1}{S_{\mathrm{tot}}(\rho)} \int_\Omega \|x - \bar{x}_{s}\|^2 \, s(x) \dif^d x.
\end{equation}
\end{definition}
The density $s$ is unrelated to the ``entanglement contour'' of Chen and Vidal \cite{chen2014entanglement}, which distributes the entanglement entropy of a region over its sites. We prove no properties of $s$ beyond those inherited from Theorem~\ref{thm:qfi_volume}.

% =========================================================================
\section{Summary}
\label{sec:synthesis}
% =========================================================================

Table~\ref{tab:master_taxonomy} lists the invariants of this volume, what is established about each, and where.

\begin{table}[htbp]
\centering
\small
\renewcommand{\arraystretch}{1.2}
\begin{tabularx}{\textwidth}{@{} l >{\raggedright\arraybackslash}p{3.4cm} >{\raggedright\arraybackslash}X @{}}
\toprule
\textbf{Invariant} & \textbf{Definition} & \textbf{What is established} \\
\midrule
$\mathcal{W}_2(A, B)$ & $\inf_\pi \int \|x-y\|^2 \dif \pi$ & Pseudometric for any choice of density realizations; sharp $L^1$ bound; metric-speed bound for the kernel realization (Thm.~\ref{thm:wasserstein_properties}). \\
\midrule
$\kappa_{\mathrm{BE}}(A)$ & $\inf \langle (-\nabla^2 \log \Phi) \xi, \xi \rangle$ & Entropy $K$-convexity (proved); log-Sobolev and variance decay (classical) (Thm.~\ref{thm:bakry_emery}). \\
\midrule
$\mathrm{Dgm}_k$, $E^{(k)}_{\mathrm{pers}}$ & Super-level persistence; entropy of bar lengths & Bottleneck stability for continuous tame realizations (classical); explicit isospectral pairs separated by the step diagrams, one also by entropy; entropy does not separate in general (Thm.~\ref{thm:isospectral_separation}, Rem.~\ref{rem:entropy_limits}). \\
\midrule
$\mathcal{W}(\Phi(A))$ & $\int \int_{\Sigma_t} H_t^2 \dif \mathcal{H}^{d-1} \dif t$ & Finite for Morse realizations, $d \ge 2$; $c_*\epsilon k^3$ law for a trigonometric checkerboard (Prop.~\ref{prop:willmore_scaling}). \\
\midrule
$h(A)$ & Weighted Cheeger constant & $\lambda_1 \ge h^2/4$ (classical, Prop.~\ref{prop:cheeger}). \\
\midrule
$\WF(\Phi(A))$ & H\"ormander wavefront set & Conormal bundle of a smooth jump interface (classical, Prop.~\ref{thm:wavefront_interface}); all directions at grid vertices of the model corner. \\
\midrule
$s^*_p(A)$ & Besov index & $1$ for $C^1$ realizations, $1/p$ for non-constant step realizations, $p > 1$ (Prop.~\ref{prop:besov}). \\
\midrule
$\Tr_\omega$, $\tau_2$, $\tau_2^\gamma$ & Dixmier trace, cocycles on $\mathbb{T}^d$ & $c_d\int\Phi$ (Connes' trace theorem); $\tau_2$ not topological; graded antisymmetrization $= -4i\deg F$ (Prop.~\ref{prop:cocycle}). \\
\midrule
$\mathrm{Vol}_{\mathrm{QFI}}(\rho)$ & $\int_\Omega \sqrt{\det g^{\mathrm{QFI}}} \dif^d x$ & Vanishes iff $\rank\dif\rho < d$ everywhere (Thm.~\ref{thm:qfi_volume}). \\
\bottomrule
\end{tabularx}
\caption{Invariants of functional realizations treated in this volume.}
\label{tab:master_taxonomy}
\end{table}

\subsection*{Open questions}
\begin{enumerate}[label=(\roman*)]
    \item Which isospectral pairs are separated by persistence of a \emph{continuous} realization, for which Theorem~\ref{thm:bottleneck_stability} applies?
    \item Upper bounds matching the lower bound $s^*_1 \ge d - \overline{\dim}_M\partial E$ for indicator realizations of rough sets.
    \item Analogues of Proposition~\ref{prop:cocycle} in higher even dimensions.
\end{enumerate}
Applications of these invariants, for example as regularizers of weight matrices or as diagnostics of tensor-network states, are not studied here.

\section*{Declarations}

\noindent\textbf{Funding.} This research was financed in part by the
Coordena\c{c}\~ao de Aperfei\c{c}oamento de Pessoal de N\'ivel Superior -- Brasil
(CAPES) -- Finance Code 001.

\medskip\noindent\textbf{Competing interests.} The author declares no competing
interests.

\medskip\noindent\textbf{Use of generative AI tools.} Large language model
assistants (Anthropic Claude; Google Gemini, through the Antigravity agent
environment) were used extensively. Under the author's direction they drafted and
edited text and \LaTeX{} source, proposed and wrote derivations, proofs and
counterexamples, wrote the numerical check scripts, and searched the literature.
The mathematics was then checked in three steps by AI sessions that had not
written it: a blind review of every statement, a correction made by a separate
session, and a targeted re-check of each correction. Every numerical claim is
checked by a script that compares it with an independent computation and includes
a negative control, a deliberately altered formula that must fail. Every reference
was resolved through Crossref, DataCite or arXiv, or, for the one book without
a DOI, through its ISBN record at Open Library. These checks were carried out by
AI systems; they are not peer review, and no result in this volume has yet
been checked by an independent human expert. AI tools are not authors. The author
chose the questions, directed the work, decided what to publish, and is
responsible for the content, including any remaining errors.

\medskip\noindent\textbf{Verification status.} Results labelled Theorem,
Proposition or Lemma are proved in the text or cited as classical results;
statements labelled Conjecture are not proved. The numerical scripts are checks,
not proofs. The \textsc{Lean 4} files in the author's repository associated with
this work are a placeholder skeleton without Mathlib and verify none of the
statements.

\medskip\noindent\textbf{Code availability.} The numerical check scripts are
available from the author on request.

% =========================================================================
\begin{thebibliography}{99}
% =========================================================================

\bibitem{ambrosio2003direct}
L.~Ambrosio and S.~Masnou,
\emph{A direct variational approach to a problem arising in image reconstruction},
Interfaces and Free Boundaries, \textbf{5} (2003), no.~1, 63--81, DOI: \href{https://doi.org/10.4171/IFB/72}{10.4171/IFB/72}.

\bibitem{ambrosio2008gradient}
L.~Ambrosio, N.~Gigli, and G.~Savar{\'e},
\emph{Gradient Flows in Metric Spaces and in the Space of Probability Measures},
2nd ed., Lectures in Mathematics ETH Z{\"u}rich, Birkh{\"a}user, Basel, 2008, DOI: \href{https://doi.org/10.1007/978-3-7643-8722-8}{10.1007/978-3-7643-8722-8}.

\bibitem{atienza2020stability}
N.~Atienza, R.~Gonzalez-D{\'\i}az, and M.~Soriano-Trigueros,
\emph{On the stability of persistent entropy and new summary functions for topological data analysis},
Pattern Recognition, \textbf{107} (2020), 107509, DOI: \href{https://doi.org/10.1016/j.patcog.2020.107509}{10.1016/j.patcog.2020.107509}.

\bibitem{bakry1985diffusions}
D.~Bakry and M.~{\'E}mery,
\emph{Diffusions hypercontractives},
S{\'e}minaire de Probabilit{\'e}s XIX 1983/84, Lecture Notes in Mathematics, vol.~1123, Springer, Berlin, Heidelberg, 1985, pp.~177--206, ISBN: 978-3-540-15230-9, DOI: \href{https://doi.org/10.1007/BFb0075847}{10.1007/BFb0075847}.

\bibitem{bakry2014analysis}
D.~Bakry, I.~Gentil, and M.~Ledoux,
\emph{Analysis and Geometry of Markov Diffusion Operators},
Grundlehren der mathematischen Wissenschaften, vol.~348, Springer, Cham, 2014, DOI: \href{https://doi.org/10.1007/978-3-319-00227-9}{10.1007/978-3-319-00227-9}.

\bibitem{benamou2000computational}
J.-D.~Benamou and Y.~Brenier,
\emph{A computational fluid mechanics solution to the Monge-Kantorovich mass transfer problem},
Numerische Mathematik, \textbf{84} (2000), no.~3, 375--393, DOI: \href{https://doi.org/10.1007/s002110050002}{10.1007/s002110050002}.

\bibitem{bhatia2019bures}
R.~Bhatia, T.~Jain, and Y.~Lim,
\emph{On the Bures--Wasserstein distance between positive definite matrices},
Expositiones Mathematicae, \textbf{37} (2019), no.~2, 165--191, DOI: \href{https://doi.org/10.1016/j.exmath.2018.01.002}{10.1016/j.exmath.2018.01.002}.

\bibitem{braunstein1994statistical}
S.~L.~Braunstein and C.~M.~Caves,
\emph{Statistical distance and the geometry of quantum states},
Physical Review Letters, \textbf{72} (1994), no.~22, 3439--3443, DOI: \href{https://doi.org/10.1103/PhysRevLett.72.3439}{10.1103/PhysRevLett.72.3439}.

\bibitem{brenier1991polar}
Y.~Brenier,
\emph{Polar factorization and monotone rearrangement of vector-valued functions},
Communications on Pure and Applied Mathematics, \textbf{44} (1991), no.~4, 375--417, DOI: \href{https://doi.org/10.1002/cpa.3160440402}{10.1002/cpa.3160440402}.

\bibitem{buser1982note}
P.~Buser,
\emph{A note on the isoperimetric constant},
Annales scientifiques de l'{\'E}cole Normale Sup{\'e}rieure (4), \textbf{15} (1982), no.~2, 213--230, DOI: \href{https://doi.org/10.24033/asens.1426}{10.24033/asens.1426}.

\bibitem{caffarelli2000monotonicity}
L.~A.~Caffarelli,
\emph{Monotonicity properties of optimal transportation and the FKG and related inequalities},
Communications in Mathematical Physics, \textbf{214} (2000), 547--563, DOI: \href{https://doi.org/10.1007/s002200000257}{10.1007/s002200000257}.

\bibitem{chan2003euler}
T.~F.~Chan, S.~H.~Kang, and J.~Shen,
\emph{Euler's elastica and curvature-based inpainting},
SIAM Journal on Applied Mathematics, \textbf{63} (2003), no.~2, 564--592, DOI: \href{https://doi.org/10.1137/S0036139901390088}{10.1137/S0036139901390088}.

\bibitem{chazal2009proximity}
F.~Chazal, D.~Cohen-Steiner, M.~Glisse, L.~J.~Guibas, and S.~Y.~Oudot,
\emph{Proximity of persistence modules and their diagrams},
in: Proceedings of the 25th Annual Symposium on Computational Geometry (SCG '09), ACM, New York, 2009, pp.~237--246, DOI: \href{https://doi.org/10.1145/1542362.1542407}{10.1145/1542362.1542407}.

\bibitem{chazal2016structure}
F.~Chazal, V.~de Silva, M.~Glisse, and S.~Oudot,
\emph{The Structure and Stability of Persistence Modules},
SpringerBriefs in Mathematics, Springer, Cham, 2016, ISBN: 978-3-319-42545-0, DOI: \href{https://doi.org/10.1007/978-3-319-42545-0}{10.1007/978-3-319-42545-0}.

\bibitem{cheeger1970lower}
J.~Cheeger,
\emph{A lower bound for the smallest eigenvalue of the Laplacian},
in: Problems in Analysis (a symposium in honor of S.~Bochner), Princeton University Press, Princeton, NJ, 1970, pp.~195--199, DOI: \href{https://doi.org/10.1515/9781400869312-013}{10.1515/9781400869312-013}.

\bibitem{chen2014entanglement}
Y.~Chen and G.~Vidal,
\emph{Entanglement contour},
Journal of Statistical Mechanics: Theory and Experiment, \textbf{2014} (2014), P10011, DOI: \href{https://doi.org/10.1088/1742-5468/2014/10/P10011}{10.1088/1742-5468/2014/10/P10011}.

\bibitem{chintakunta2015entropy}
H.~Chintakunta, T.~Gentimis, R.~Gonzalez-D{\'\i}az, M.-J.~Jim{\'e}nez, and H.~Krim,
\emph{An entropy-based persistence barcode},
Pattern Recognition, \textbf{48} (2015), no.~2, 391--401, DOI: \href{https://doi.org/10.1016/j.patcog.2014.06.023}{10.1016/j.patcog.2014.06.023}.

\bibitem{chowdhury2019gromov}
S.~Chowdhury and F.~M{\'e}moli,
\emph{The Gromov--Wasserstein distance between networks and stable network invariants},
Information and Inference: A Journal of the IMA, \textbf{8} (2019), no.~4, 757--787, DOI: \href{https://doi.org/10.1093/imaiai/iaz026}{10.1093/imaiai/iaz026}.

\bibitem{cohensteiner2007stability}
D.~Cohen-Steiner, H.~Edelsbrunner, and J.~Harer,
\emph{Stability of persistence diagrams},
Discrete \& Computational Geometry, \textbf{37} (2007), 103--120, DOI: \href{https://doi.org/10.1007/s00454-006-1276-5}{10.1007/s00454-006-1276-5}.

\bibitem{connes1985noncommutative}
A.~Connes,
\emph{Non-commutative differential geometry},
Publications Math{\'e}matiques de l'IH{\'E}S, \textbf{62} (1985), 41--144, DOI: \href{https://doi.org/10.1007/BF02698807}{10.1007/BF02698807}.

\bibitem{connes1988action}
A.~Connes,
\emph{The action functional in non-commutative geometry},
Communications in Mathematical Physics, \textbf{117} (1988), 673--683, DOI: \href{https://doi.org/10.1007/BF01218391}{10.1007/BF01218391}.

\bibitem{connes1994noncommutative}
A.~Connes,
\emph{Noncommutative Geometry},
Academic Press, San Diego, CA, 1994, ISBN: 978-0-12-185860-5.

\bibitem{connes1995local}
A.~Connes and H.~Moscovici,
\emph{The local index formula in noncommutative geometry},
Geometric and Functional Analysis, \textbf{5} (1995), no.~2, 174--243, DOI: \href{https://doi.org/10.1007/BF01895667}{10.1007/BF01895667}.

\bibitem{crawley2015decomposition}
W.~Crawley-Boevey,
\emph{Decomposition of pointwise finite-dimensional persistence modules},
Journal of Algebra and Its Applications, \textbf{14} (2015), 1550066, DOI: \href{https://doi.org/10.1142/S0219498815500668}{10.1142/S0219498815500668}.

\bibitem{edelsbrunner2002topological}
H.~Edelsbrunner, D.~Letscher, and A.~Zomorodian,
\emph{Topological persistence and simplification},
Discrete \& Computational Geometry, \textbf{28} (2002), no.~4, 511--533, DOI: \href{https://doi.org/10.1007/s00454-002-2885-2}{10.1007/s00454-002-2885-2}.

\bibitem{edelsbrunner2010computational}
H.~Edelsbrunner and J.~L.~Harer,
\emph{Computational Topology: An Introduction},
American Mathematical Society, Providence, RI, 2010, ISBN: 978-0-8218-4925-5, DOI: \href{https://doi.org/10.1090/mbk/069}{10.1090/mbk/069}.

\bibitem{evans2015measure}
L.~C.~Evans and R.~F.~Gariepy,
\emph{Measure Theory and Fine Properties of Functions},
revised ed., CRC Press, Boca Raton, FL, 2015, DOI: \href{https://doi.org/10.1201/b18333}{10.1201/b18333}.

\bibitem{gordon1992one}
C.~Gordon, D.~L.~Webb, and S.~Wolpert,
\emph{One cannot hear the shape of a drum},
Bulletin of the American Mathematical Society (N.S.), \textbf{27} (1992), no.~1, 134--138, DOI: \href{https://doi.org/10.1090/S0273-0979-1992-00289-6}{10.1090/S0273-0979-1992-00289-6}.

\bibitem{grochenig2001foundations}
K.~Gr{\"o}chenig,
\emph{Foundations of Time-Frequency Analysis},
Applied and Numerical Harmonic Analysis, Birkh{\"a}user, Boston, MA, 2001, ISBN: 978-1-4612-6568-9, DOI: \href{https://doi.org/10.1007/978-1-4612-0003-1}{10.1007/978-1-4612-0003-1}.

\bibitem{guillemin1985new}
V.~Guillemin,
\emph{A new proof of Weyl's formula on the asymptotic distribution of eigenvalues},
Advances in Mathematics, \textbf{55} (1985), no.~2, 131--160, DOI: \href{https://doi.org/10.1016/0001-8708(85)90018-0}{10.1016/0001-8708(85)90018-0}.

\bibitem{helstrom1967minimum}
C.~W.~Helstrom,
\emph{Minimum mean-squared error of estimates in quantum statistics},
Physics Letters A, \textbf{25} (1967), no.~2, 101--102, DOI: \href{https://doi.org/10.1016/0375-9601(67)90366-0}{10.1016/0375-9601(67)90366-0}.

\bibitem{hormander2003analysis}
L.~H{\"o}rmander,
\emph{The Analysis of Linear Partial Differential Operators I: Distribution Theory and Fourier Analysis},
Classics in Mathematics, Springer-Verlag, Berlin, Heidelberg, 2003, ISBN: 978-3-540-00662-6, DOI: \href{https://doi.org/10.1007/978-3-642-61497-2}{10.1007/978-3-642-61497-2}.

\bibitem{jordan1998variational}
R.~Jordan, D.~Kinderlehrer, and F.~Otto,
\emph{The variational formulation of the Fokker--Planck equation},
SIAM Journal on Mathematical Analysis, \textbf{29} (1998), no.~1, 1--17, DOI: \href{https://doi.org/10.1137/S0036141096303359}{10.1137/S0036141096303359}.

\bibitem{kac1966can}
M.~Kac,
\emph{Can one hear the shape of a drum?},
The American Mathematical Monthly, \textbf{73} (1966), no.~4, part~II, 1--23, DOI: \href{https://doi.org/10.1080/00029890.1966.11970915}{10.1080/00029890.1966.11970915}.

\bibitem{kaczynski2004computational}
T.~Kaczynski, K.~Mischaikow, and M.~Mrozek,
\emph{Computational Homology},
Applied Mathematical Sciences, Springer, New York, 2004, ISBN: 978-0-387-21597-6, DOI: \href{https://doi.org/10.1007/b97315}{10.1007/b97315}.

\bibitem{kantorovich2006translocation}
L.~V.~Kantorovich,
\emph{On the translocation of masses},
Journal of Mathematical Sciences, \textbf{133} (2006), no.~4, 1381--1382, DOI: \href{https://doi.org/10.1007/s10958-006-0049-2}{10.1007/s10958-006-0049-2}. English translation of Dokl.\ Akad.\ Nauk SSSR \textbf{37} (1942), 227--229.

\bibitem{ledoux1994simple}
M.~Ledoux,
\emph{A simple analytic proof of an inequality by P.~Buser},
Proceedings of the American Mathematical Society, \textbf{121} (1994), no.~3, 951--959, DOI: \href{https://doi.org/10.1090/S0002-9939-1994-1186991-X}{10.1090/S0002-9939-1994-1186991-X}.

\bibitem{liu2020quantum}
J.~Liu, H.~Yuan, X.-M.~Lu, and X.~Wang,
\emph{Quantum Fisher information matrix and multiparameter estimation},
Journal of Physics A: Mathematical and Theoretical, \textbf{53} (2020), no.~2, 023001, DOI: \href{https://doi.org/10.1088/1751-8121/ab5d4d}{10.1088/1751-8121/ab5d4d}, arXiv: \href{https://arxiv.org/abs/1907.08037}{1907.08037}.

\bibitem{lombardini2019fractional}
L.~Lombardini,
\emph{Fractional perimeters from a fractal perspective},
Advanced Nonlinear Studies, \textbf{19} (2019), no.~1, 165--196, DOI: \href{https://doi.org/10.1515/ans-2018-2016}{10.1515/ans-2018-2016}, arXiv: \href{https://arxiv.org/abs/1603.06088}{1603.06088}.

\bibitem{lord2013singular}
S.~Lord, F.~Sukochev, and D.~Zanin,
\emph{Singular Traces: Theory and Applications},
De Gruyter Studies in Mathematics, De Gruyter, Berlin, 2012, ISBN: 978-3-11-026250-6, DOI: \href{https://doi.org/10.1515/9783110262551}{10.1515/9783110262551}.

\bibitem{lott2009ricci}
J.~Lott and C.~Villani,
\emph{Ricci curvature for metric-measure spaces via optimal transport},
Annals of Mathematics, \textbf{169} (2009), no.~3, 903--991, DOI: \href{https://doi.org/10.4007/annals.2009.169.903}{10.4007/annals.2009.169.903}.

\bibitem{lovasz2012large}
L.~Lov{\'a}sz,
\emph{Large Networks and Graph Limits},
Colloquium Publications, vol.~60, American Mathematical Society, Providence, RI, 2012, DOI: \href{https://doi.org/10.1090/coll/060}{10.1090/coll/060}.

\bibitem{marques2014minmax}
F.~C.~Marques and A.~Neves,
\emph{Min-max theory and the Willmore conjecture},
Annals of Mathematics, \textbf{179} (2014), no.~2, 683--782, DOI: \href{https://doi.org/10.4007/annals.2014.179.2.6}{10.4007/annals.2014.179.2.6}.

\bibitem{mccann1997convexity}
R.~J.~McCann,
\emph{A convexity principle for interacting gases},
Advances in Mathematics, \textbf{128} (1997), 153--179, DOI: \href{https://doi.org/10.1006/aima.1997.1634}{10.1006/aima.1997.1634}.

\bibitem{memoli2011gromov}
F.~M{\'e}moli,
\emph{Gromov--Wasserstein distances and the metric approach to object matching},
Foundations of Computational Mathematics, \textbf{11} (2011), 417--487, DOI: \href{https://doi.org/10.1007/s10208-011-9093-5}{10.1007/s10208-011-9093-5}.

\bibitem{otto2000generalization}
F.~Otto and C.~Villani,
\emph{Generalization of an inequality by Talagrand and links with the logarithmic Sobolev inequality},
Journal of Functional Analysis, \textbf{173} (2000), no.~2, 361--400, DOI: \href{https://doi.org/10.1006/jfan.1999.3557}{10.1006/jfan.1999.3557}.

\bibitem{petz1996monotone}
D.~Petz,
\emph{Monotone metrics on matrix spaces},
Linear Algebra and its Applications, \textbf{244} (1996), 81--96, DOI: \href{https://doi.org/10.1016/0024-3795(94)00211-8}{10.1016/0024-3795(94)00211-8}.

\bibitem{peyre2018comparison}
R.~Peyre,
\emph{Comparison between $W_2$ distance and $\dot{H}^{-1}$ norm, and localization of Wasserstein distance},
ESAIM: Control, Optimisation and Calculus of Variations, \textbf{24} (2018), no.~4, 1489--1501, DOI: \href{https://doi.org/10.1051/cocv/2017050}{10.1051/cocv/2017050}, arXiv: \href{https://arxiv.org/abs/1104.4631}{1104.4631}.

\bibitem{provost1980riemannian}
J.~P.~Provost and G.~Vallee,
\emph{Riemannian structure on manifolds of quantum states},
Communications in Mathematical Physics, \textbf{76} (1980), no.~3, 289--301, DOI: \href{https://doi.org/10.1007/BF02193559}{10.1007/BF02193559}.

\bibitem{safranek2017discontinuities}
D.~{\v{S}}afr{\'a}nek,
\emph{Discontinuities of the quantum Fisher information and the Bures metric},
Physical Review A, \textbf{95} (2017), no.~5, 052320, DOI: \href{https://doi.org/10.1103/PhysRevA.95.052320}{10.1103/PhysRevA.95.052320}.

\bibitem{sard1942measure}
A.~Sard,
\emph{The measure of the critical values of differentiable maps},
Bulletin of the American Mathematical Society, \textbf{48} (1942), no.~12, 883--890, DOI: \href{https://doi.org/10.1090/S0002-9904-1942-07811-6}{10.1090/S0002-9904-1942-07811-6}.

\bibitem{silvafilho2026beyond}
R.~M.~Silva-Filho,
\emph{Beyond the Spectrum: Functional Realizations of Matrices and Tensors, Emerging Invariants, and Geometric Measures},
Volume~I of \emph{Beyond the Spectrum}, Zenodo \href{https://doi.org/10.5281/zenodo.22644743}{10.5281/zenodo.22644743}, 2026; this Zenodo record holds all three volumes of the series.

\bibitem{silvafilho2026geometry}
R.~M.~Silva-Filho,
\emph{Geometry, Tensors, and Quantum Gravity},
monograph, Zenodo, 2026, DOI: \href{https://doi.org/10.5281/zenodo.22290043}{10.5281/zenodo.22290043}.

\bibitem{sturm2006geometryI}
K.-T.~Sturm,
\emph{On the geometry of metric measure spaces. I},
Acta Mathematica, \textbf{196} (2006), no.~1, 65--131, DOI: \href{https://doi.org/10.1007/s11511-006-0002-8}{10.1007/s11511-006-0002-8}.

\bibitem{sturm2006geometryII}
K.-T.~Sturm,
\emph{On the geometry of metric measure spaces. II},
Acta Mathematica, \textbf{196} (2006), no.~1, 133--177, DOI: \href{https://doi.org/10.1007/s11511-006-0003-7}{10.1007/s11511-006-0003-7}.

\bibitem{sturm2023space}
K.-T.~Sturm,
\emph{The Space of Spaces: Curvature Bounds and Gradient Flows on the Space of Metric Measure Spaces},
Memoirs of the American Mathematical Society, \textbf{290} (2023), no.~1443, DOI: \href{https://doi.org/10.1090/memo/1443}{10.1090/memo/1443}.

\bibitem{triebel1983theory}
H.~Triebel,
\emph{Theory of Function Spaces},
Monographs in Mathematics, Birkh{\"a}user, Basel, 1983, ISBN: 978-3-0346-0415-4, DOI: \href{https://doi.org/10.1007/978-3-0346-0416-1}{10.1007/978-3-0346-0416-1}.

\bibitem{uhlmann1976transition}
A.~Uhlmann,
\emph{The ``transition probability'' in the state space of a $*$-algebra},
Reports on Mathematical Physics, \textbf{9} (1976), no.~2, 273--279, DOI: \href{https://doi.org/10.1016/0034-4877(76)90060-4}{10.1016/0034-4877(76)90060-4}.

\bibitem{vandam2003which}
E.~R.~van Dam and W.~H.~Haemers,
\emph{Which graphs are determined by their spectrum?},
Linear Algebra and its Applications, \textbf{373} (2003), 241--272, DOI: \href{https://doi.org/10.1016/S0024-3795(03)00483-X}{10.1016/S0024-3795(03)00483-X}.

\bibitem{verstraete2010continuous}
F.~Verstraete and J.~I.~Cirac,
\emph{Continuous matrix product states for quantum fields},
Physical Review Letters, \textbf{104} (2010), no.~19, 190405, DOI: \href{https://doi.org/10.1103/PhysRevLett.104.190405}{10.1103/PhysRevLett.104.190405}.

\bibitem{villani2009optimal}
C.~Villani,
\emph{Optimal Transport: Old and New},
Grundlehren der mathematischen Wissenschaften, vol.~338, Springer-Verlag, Berlin, Heidelberg, 2009, ISBN: 978-3-540-71049-3, DOI: \href{https://doi.org/10.1007/978-3-540-71050-9}{10.1007/978-3-540-71050-9}.

\bibitem{visintin1991generalized}
A.~Visintin,
\emph{Generalized coarea formula and fractal sets},
Japan Journal of Industrial and Applied Mathematics, \textbf{8} (1991), no.~2, 175--201, DOI: \href{https://doi.org/10.1007/BF03167679}{10.1007/BF03167679}.

\bibitem{von2005transport}
M.-K.~von Renesse and K.-T.~Sturm,
\emph{Transport inequalities, gradient estimates, entropy and Ricci curvature},
Communications on Pure and Applied Mathematics, \textbf{58} (2005), no.~7, 923--940, DOI: \href{https://doi.org/10.1002/cpa.20060}{10.1002/cpa.20060}.

\bibitem{wagner2012efficient}
H.~Wagner, C.~Chen, and E.~Vu{\c{c}}ini,
\emph{Efficient computation of persistent homology for cubical data},
in: Topological Methods in Data Analysis and Visualization II, Mathematics and Visualization, Springer, Berlin, Heidelberg, 2012, pp.~91--106, DOI: \href{https://doi.org/10.1007/978-3-642-23175-9_7}{10.1007/978-3-642-23175-9\_7}.

\bibitem{white1973global}
J.~H.~White,
\emph{A global invariant of conformal mappings in space},
Proceedings of the American Mathematical Society, \textbf{38} (1973), 162--164, DOI: \href{https://doi.org/10.1090/S0002-9939-1973-0324603-1}{10.1090/S0002-9939-1973-0324603-1}.

\bibitem{wodzicki1987noncommutative}
M.~Wodzicki,
\emph{Noncommutative residue. Chapter I. Fundamentals},
in: K-Theory, Arithmetic and Geometry (Moscow, 1984--1986), Lecture Notes in Mathematics, vol.~1289, Springer, Berlin, 1987, pp.~320--399, DOI: \href{https://doi.org/10.1007/BFb0078372}{10.1007/BFb0078372}.

\bibitem{zanardi2007information}
P.~Zanardi, P.~Giorda, and M.~Cozzini,
\emph{Information-theoretic differential geometry of quantum phase transitions},
Physical Review Letters, \textbf{99} (2007), no.~10, 100603, DOI: \href{https://doi.org/10.1103/PhysRevLett.99.100603}{10.1103/PhysRevLett.99.100603}.

\bibitem{zomorodian2005computing}
A.~Zomorodian and G.~Carlsson,
\emph{Computing persistent homology},
Discrete \& Computational Geometry, \textbf{33} (2005), no.~2, 249--274, DOI: \href{https://doi.org/10.1007/s00454-004-1146-y}{10.1007/s00454-004-1146-y}.

\end{thebibliography}

\end{document}
